% ============================================================================
% AxicCReview.tex --- Articulo IEEE, Etapa 1 (preliminar)
%
% VERSION INDIVIDUAL --- Sebastian Hernandez Bonilla, Eje C.
% Construida sobre los siete trabajos arbitrados asignados a este autor.
% Es la entrada para la fusion a tres bandas que producira el articulo
% conjunto de la Etapa 1; no es ese articulo.
%
% DOCUMENTO AUTOCONTENIDO. Todo el contenido academico vive en este
% fichero: preambulo, secciones, tablas y figuras. La unica dependencia
% externa es la bibliografia.
%
% Compilacion:  pdflatex -> bibtex -> pdflatex -> pdflatex
% Ficheros necesarios: AxicCReview.tex + references.bib
% ============================================================================
\documentclass[conference]{IEEEtran}


\usepackage[T1]{fontenc}
\usepackage[utf8]{inputenc}
\usepackage[spanish,es-tabla]{babel}
% Conserva el punto como separador decimal, para que las cifras coincidan
% exactamente con las reportadas por las fuentes citadas y por nuestra
% auditoria. Si babel-spanish no define la orden, la linea es inocua.
\providecommand{\decimalpoint}{}
\decimalpoint
\usepackage{microtype}


\usepackage{amsmath}
\usepackage{amssymb}


\usepackage{array}
\usepackage{booktabs}
\usepackage{array}
\usepackage{tabularx}
\usepackage{enumitem}

\newcolumntype{Y}{>{\raggedright\arraybackslash}X}
\newcolumntype{P}[1]{>{\raggedright\arraybackslash}p{#1}}
\renewcommand{\arraystretch}{1.15}
\setlist[itemize]{leftmargin=*,noitemsep,topsep=2pt}
\setlist[enumerate]{leftmargin=*,noitemsep,topsep=2pt}


\usepackage{xcolor}
\usepackage{cite}
\usepackage[hidelinks]{hyperref}


\newcommand{\code}[1]{\texttt{#1}}
\newcommand{\pendiente}[1]{%
  \texorpdfstring{\textcolor{red}{\textbf{[POR COMPLETAR: #1]}}}%
                 {[POR COMPLETAR: #1]}}

% Rotulo de palabras clave en espanol (IEEEtran lo fija en ingles).
\providecommand{\IEEEkeywordsname}{Index Terms}
\renewcommand{\IEEEkeywordsname}{Palabras clave}

\begin{document}

\title{Estimaci\'on de la probabilidad de victoria a partir del estado de
juego en el \emph{Pok\'emon Trading Card Game}: un estudio tabular con
control de fuga de datos sobre repeticiones de agentes competitivos\\
{\footnotesize \textit{Versi\'on preliminar --- Etapa 1: estado del arte y
dise\~no propuesto. Revisi\'on individual del Eje C (conjunto de datos y
validaci\'on).}}}

\author{
\IEEEauthorblockN{Sebasti\'an Hern\'andez Bonilla}
\IEEEauthorblockA{\textit{Escuela de Ingenier\'ia en Computaci\'on} \\
\textit{Instituto Tecnol\'ogico de Costa Rica}\\
Cartago, Costa Rica \\
2022093651}
}

\maketitle

% ===========================================================================
% 6.2.1  RESUMEN PRELIMINAR Y PALABRAS CLAVE
% ===========================================================================
\begin{abstract}
\textbf{Resumen preliminar.} Este resumen describe el problema, la
motivaci\'on y el enfoque propuesto. No reporta resultados experimentales,
puesto que en esta etapa no se ha entrenado ning\'un modelo.
Evaluar una posici\'on es una competencia central en los juegos de cartas
competitivos y, sin embargo, en el \emph{Pok\'emon Trading Card Game}
(PTCG) no existe trabajo publicado que estime la probabilidad de victoria
a partir de un estado de juego intermedio y parcialmente observable. El
corpus de repeticiones del PTCG AI Battle Challenge, publicado
recientemente, vuelve tratable la pregunta, pero su estructura hace que
sea f\'acil responderla mal: una sola partida genera miles de estados
fuertemente correlacionados, de modo que una partici\'on aleatoria a nivel
de fila coloca estados de la misma partida a ambos lados de la divisi\'on.
En una auditor\'ia de 52\,085 episodios medimos que tal partici\'on
situar\'ia el 99.995\% de los episodios en ambas particiones de forma
simult\'anea. Este trabajo pregunta, por tanto, cu\'anta informaci\'on
sobre el desenlace contiene un \'unico estado observable, c\'omo crece esa
informaci\'on conforme avanza la partida y cu\'anta sobrevive al evaluar
sobre partidas, mazos y periodos no vistos. Proponemos una formulaci\'on
tabular supervisada, una fila por estado fresco del asiento activo, con
caracter\'isticas derivadas exclusivamente de la observaci\'on disponible
en ese instante, con una escalera de modelos cl\'asicos que va desde un
predictor trivial hasta \emph{gradient boosting}, una lista negra
expl\'icita de campos cuyo uso constituir\'ia fuga de datos, y particiones
agrupadas por episodio en lugar de por fila. La evaluaci\'on reporta
discriminaci\'on y calibraci\'on por separado, de forma global y por
ventana de turno, con intervalos de confianza remuestreados sobre
episodios. El \emph{baseline} de literatura es la tarea an\'aloga sobre
estados de juego de \emph{Hearthstone}, cuyo \emph{baseline} oficial
alcanz\'o un AUC de 0.7846 y cuya mejor soluci\'on alcanz\'o 0.8019; las
diferencias de juego, protocolo y poblaci\'on de agentes que impiden una
comparaci\'on estricta se declaran expl\'icitamente.
\end{abstract}

\begin{IEEEkeywords}
estimaci\'on de probabilidad de victoria, juegos de cartas coleccionables,
aprendizaje autom\'atico tabular, fuga de datos, validaci\'on cruzada
agrupada, cambio de distribuci\'on
\end{IEEEkeywords}

% ===========================================================================
% 6.2.2  INTRODUCCION
% ===========================================================================
\section{Introducci\'on}
\label{sec:introduction}

\subsection{Contexto y relevancia}
\label{subsec:context}

Estimar la probabilidad de que un jugador gane desde la posici\'on que
tiene actualmente sobre la mesa es una de las competencias recurrentes de
los sistemas que juegan. Una funci\'on de evaluaci\'on de estados convierte
una b\'usqueda intratable en una tratable, y es lo que permite a un agente
detener una simulaci\'on antes de tiempo y a\'un as\'i devolver un juicio
con sentido. En los juegos de cartas competitivos esa misma estimaci\'on
tiene un segundo p\'ublico: analistas, comentaristas y jugadores la leen
como una afirmaci\'on sobre qu\'e recursos deciden realmente una partida.
Ambos usos dependen del mismo objeto, una funci\'on que asigna a una
posici\'on observable una probabilidad de victoria.

Dos propiedades convierten a los juegos de cartas coleccionables (CCG, por
sus siglas en ingl\'es) en un escenario dif\'icil para esa funci\'on. La
primera es la informaci\'on imperfecta: un jugador nunca ve la mano del
rival ni las cartas de premio boca abajo, de modo que el evaluador debe
trabajar sobre una vista parcial y no sobre el estado verdadero del mundo.
La segunda es la no estacionariedad. El repertorio de cartas legales, los
mazos que los jugadores efectivamente llevan y la poblaci\'on de
adversarios cambian con el tiempo, y un modelo ajustado a un metajuego no
tiene por qu\'e sobrevivir al siguiente. Bertram \textit{et al.} lo
muestran de forma concreta para \emph{Magic: The Gathering}: un modelo
cuya representaci\'on de carta es una codificaci\'on \emph{one-hot}
alcanza un 67.80\% de acierto \emph{top-1} sobre el conjunto de cartas
con el que fue entrenado y queda simplemente indefinido sobre cartas que
nunca ha visto, mientras que una representaci\'on generalizada conserva un
$68.00\%$ dentro de la distribuci\'on y a\'un alcanza un $42.87\%$
sobre conjuntos no vistos~\cite{bertram_cardrepresentations_2024}.

El problema de la evaluaci\'on ya se ha planteado antes en esta familia de
juegos. El AAIA'17 Data Mining Challenge pidi\'o a los participantes
predecir el ganador de una partida de \emph{Hearthstone} a partir de un
\'unico estado de juego, y report\'o un \emph{baseline} oficial de 0.7846
de AUC frente a una mejor soluci\'on de
0.8019~\cite{janusz_hearthstone_2017}. Para \emph{StarCraft II}, el
descriptor de datos SC2EGSet acompa\~na su corpus de repeticiones con
experimentos de predicci\'on de victoria sobre indicadores econ\'omicos
agregados, y alcanza un acierto de 0.9071 cuando dispone de las
caracter\'isticas de ambos jugadores~\cite{bialecki_sc2egset_2023}. El
\emph{Pok\'emon Trading Card Game} (PTCG) carece de un estudio publicado
comparable, pese a que su estructura de recursos, cartas de premio,
asignaci\'on de energ\'ia, banca, l\'ineas de evoluci\'on, difiere de la
de ambos juegos de maneras que deber\'ian cambiar qu\'e caracter\'isticas
transportan informaci\'on.

Esa brecha se ha vuelto abordable recientemente. El PTCG AI Battle
Challenge publica a diario repeticiones JSON de partidas completas entre
agentes que compiten entre s\'i~\cite{ptcg_dataset_2026}. El corpus es
grande y, para esta tarea, estructuralmente traicionero. Nuestra
auditor\'ia de diez publicaciones diarias, que cubren 52\,085 episodios,
encontr\'o que esos episodios producen 7\,038\,378 instant\'aneas de
estado utilizables: aproximadamente 135 filas correlacionadas por cada
partida independiente. Datos con esta forma premian el protocolo
equivocado. Karbalaie \textit{et al.} midieron el efecto directamente en un
conjunto cl\'inico de medidas repetidas, donde un clasificador
\emph{extra trees} obtuvo 0.91 de acierto bajo validaci\'on cruzada
estratificada de diez pliegues y 0.66 bajo validaci\'on dejando un
participante fuera, una inflaci\'on de 0.25 producida enteramente por la
regla de partici\'on~\cite{karbalaie_participantaware_2026}. El costo de
ignorar la estructura no es, por tanto, un sesgo marginal: es la diferencia
entre un resultado y un artefacto.

\subsection{Planteamiento del problema}
\label{subsec:problem}

Este trabajo ataca una limitaci\'on espec\'ifica: para el \emph{Pok\'emon
TCG} no existe evidencia sobre cu\'anta informaci\'on contiene un \'unico
estado observable acerca del desenlace, sobre c\'omo se acumula esa
informaci\'on a lo largo de la partida, ni sobre cu\'anta se transfiere a
partidas, mazos y periodos temporales que no se vieron durante el
entrenamiento. La literatura que responder\'ia esas preguntas para juegos
adyacentes responde solo una parte. Janusz \textit{et al.} reportan
discriminaci\'on pero no calibraci\'on~\cite{janusz_hearthstone_2017};
Bia\l ecki \textit{et al.} reportan acierto y describen sus propios
experimentos como ``m\'as ilustrativos que
investigativos''~\cite{bialecki_sc2egset_2023}; y ninguno reporta
rendimiento bajo un cambio de distribuci\'on declarado.

El problema es tambi\'en metodol\'ogico. Puesto que el corpus fue
recolectado por un tercero y publicado en medio de una competencia en
curso, no controlamos su procedencia, que es precisamente la situaci\'on
que Kaufman \textit{et al.} identifican como el caso m\'as dif\'icil para
el control de la fuga de datos~\cite{kaufman_leakage_2011}. Establecer que
una cifra reportada significa lo que aparenta significar forma parte, por
tanto, del problema de investigaci\'on y no de sus preliminares.

\subsection{Pregunta de investigaci\'on}
\label{subsec:rq}
\begin{quote}
\textbf{PI:} \textit{\textquestiondown En qu\'e medida las caracter\'isticas observables
del estado de una partida de \emph{Pok\'emon TCG} permiten predecir su
resultado final mediante modelos cl\'asicos de aprendizaje autom\'atico;
c\'omo cambia esa capacidad predictiva conforme avanza la partida; y
cu\'anta se conserva al evaluar sobre episodios, mazos y periodos
temporales ausentes del entrenamiento?}
\end{quote}

\noindent La pregunta admite refutaci\'on, es medible y puede responderse
dentro del semestre. Es refutable porque el modelo puede no superar a un
predictor trivial en los turnos tempranos; es medible mediante ROC-AUC y
\emph{log loss} calculados por ventana de turno sobre particiones
agrupadas por episodio; y es respondible porque los datos ya est\'an
descargados y auditados, y definir el protocolo no requiere entrenar
ning\'un modelo.

\subsection{Hip\'otesis de trabajo}
\label{subsec:hypothesis}
\begin{quote}
\textbf{H1:} Las caracter\'isticas estructuradas del estado observable
contienen se\~nal predictiva sobre el desenlace desde etapas tempranas de
la partida, ROC-AUC significativamente superior a $0.5$ en la ventana de
turnos temprana, y esa se\~nal aumenta conforme la partida progresa.

\textbf{H2:} La se\~nal temprana no se explica \'unicamente por el orden de
turno ni por la identidad del mazo.

\textbf{H3:} Bajo el protocolo temporal, la calibraci\'on se degrada antes
y en mayor medida que la discriminaci\'on.
\end{quote}

\noindent H3 es una predicci\'on importada de un campo vecino, no una
conjetura. Guo \textit{et al.} revisaron quince estudios cl\'inicos sobre
cambio temporal del conjunto de datos y encontraron que los once estudios
que midieron calibraci\'on reportaron deterioro de la calibraci\'on,
mientras que solo tres de los doce que midieron discriminaci\'on
reportaron deterioro de la
discriminaci\'on~\cite{guo_temporalshift_2021}. Si esa asimetr\'ia se
sostiene en un dominio de juegos es una pregunta emp\'irica abierta, y
enunciarla como hip\'otesis la vuelve falsable aqu\'i.

Cada hip\'otesis tiene un criterio de refutaci\'on declarado. H1 se debilita
si el intervalo de confianza del AUC de la ventana temprana incluye $0.5$,
y su componente din\'amico queda refutado si el AUC por ventana no crece
una vez controlada la duraci\'on de la partida. H2 queda refutada si un
modelo restringido al orden de turno y a la identidad del mazo iguala al
conjunto completo de caracter\'isticas, lo que indicar\'ia un artefacto de
emparejamiento y no se\~nal del estado de juego. H3 queda refutada si el
\emph{Brier score} y las curvas de fiabilidad no se degradan m\'as
r\'apido que el AUC bajo la partici\'on temporal. No se fija de antemano
ning\'un umbral num\'erico de mejora: sin resultado emp\'irico previo sobre
este corpus, cualquier n\'umero de ese tipo ser\'ia inventado.

\subsection{Objetivos}
\label{subsec:objectives}
\textbf{Objetivo general.} Desarrollar y evaluar modelos cl\'asicos de
aprendizaje autom\'atico que estimen la probabilidad de victoria en
partidas de \emph{Pok\'emon TCG} a partir de caracter\'isticas tabulares
derivadas exclusivamente del estado observable, y caracterizar c\'omo
evoluciona esa capacidad predictiva a lo largo de la partida y ante
cambios de episodio, mazo y periodo temporal.

\noindent\textbf{Objetivos espec\'ificos.}
\begin{enumerate}
  \item Caracterizar el corpus y su calidad sobre la representaci\'on
        tabular final: an\'alisis exploratorio, diccionario de datos,
        mecanismos de ausencia, valores at\'ipicos, balance de clases y
        representatividad.
  \item Construir una tuber\'ia reproducible que vaya de la repetici\'on
        JSON a una tabla de instant\'aneas frescas, con perspectiva
        normalizada, semillas fijas y verificaci\'on de esquema.
  \item Dise\~nar y verificar los controles de fuga de datos: agrupaci\'on
        por episodio, lista negra expl\'icita de campos prohibidos y
        pruebas automatizadas que fallen si una columna prohibida entra al
        vector de caracter\'isticas.
  \item Comparar una escalera de \emph{baselines}, desde un predictor
        trivial hasta \emph{gradient boosting}, bajo tres protocolos de
        partici\'on, reportando discriminaci\'on y calibraci\'on con
        intervalos de confianza calculados sobre episodios.
  \item Analizar explicabilidad y generalizaci\'on: ablaciones por familia
        de caracter\'isticas, importancia por permutaci\'on agrupada,
        curvas de fiabilidad y la curva de capacidad predictiva por turno.
\end{enumerate}

\subsection{Contribuci\'on esperada}
\label{subsec:contribution}

Se esperan tres contribuciones. Primera, una caracterizaci\'on de la
probabilidad de victoria desde el estado de juego en un juego de cartas que
no ha sido estudiado de esta forma, incluyendo el perfil por turno de la
capacidad predictiva que Janusz \textit{et al.} reportan para
\emph{Hearthstone} pero que no existe para el \emph{PTCG}. Segunda, un
protocolo de evaluaci\'on cuyos controles de fuga est\'an medidos y no
solo afirmados: la auditor\'ia cuantifica la contaminaci\'on que
producir\'ia el protocolo ingenuo, de modo que el dise\~no agrupado se
justifica con una cifra de este corpus y no por analog\'ia. Tercera, una
tuber\'ia reproducible y un diccionario de datos documentado para un corpus
que hoy carece de descripci\'on arbitrada, siguiendo la estructura que el
descriptor SC2EGSet establece para datos de repeticiones de deportes
electr\'onicos~\cite{bialecki_sc2egset_2023}.

Un resultado negativo tambi\'en ser\'ia una contribuci\'on. Si la
discriminaci\'on colapsa bajo el protocolo temporal, eso constituye
evidencia sobre la rapidez con que un metajuego invalida un evaluador de
estados, que es la pregunta que Bertram \textit{et al.} plantean para
repertorios de cartas en cambio
permanente~\cite{bertram_cardrepresentations_2024}.

\subsection{Estructura del documento}
\label{subsec:structure}
La Secci\'on~\ref{sec:related} revisa la literatura a lo largo de tres ejes
y cierra con el \emph{baseline} de literatura y las brechas identificadas.
La Secci\'on~\ref{sec:methodology} presenta la metodolog\'ia propuesta. La
Secci\'on~\ref{sec:trackanalysis} detalla el an\'alisis t\'ecnico que exige
el Track A. La Secci\'on~\ref{sec:ethics} discute las consideraciones
\'eticas preliminares. La Secci\'on~\ref{sec:aideclaration} declara el uso
de herramientas de inteligencia artificial. La
Secci\'on~\ref{sec:conclusion} cierra con el estado actual y los pr\'oximos
pasos.

% ===========================================================================
% 6.2.3  TRABAJOS RELACIONADOS --- TRES EJES
% ===========================================================================
\section{Trabajos Relacionados}
\label{sec:related}

Esta revisi\'on cubre los siete trabajos arbitrados asignados al autor bajo
el Eje C (conjunto de datos, validaci\'on y cambio de distribuci\'on), que
es el eje del que depende la validez de todo el estudio. Los candidatos se
recuperaron de Scopus, IEEE Xplore, la ACM Digital Library y los registros
de los editores de las revistas correspondientes, mediante cadenas de
b\'usqueda que combinaron \code{win prediction}, \code{game state},
\code{collectible card game} y \code{esport replay dataset} para el
dominio, y \code{data leakage}, \code{grouped cross-validation},
\code{repeated measures} y \code{temporal dataset shift} para la
metodolog\'ia. Se examinaron nueve candidatos a texto completo y se
retuvieron siete. La inclusi\'on exigi\'o venue arbitrado de revista o
conferencia, un DOI que resolviera, y contenido que respaldara una
decisi\'on efectivamente tomada en este trabajo; las dos exclusiones est\'an
documentadas en el informe de selecci\'on del autor. Cinco de los siete se
publicaron en 2021 o despu\'es. Los dos trabajos m\'as antiguos se retienen
deliberadamente: uno es la formulaci\'on can\'onica de la fuga de
datos~\cite{kaufman_leakage_2011}, el otro es la \'unica competencia
publicada de predicci\'on de victoria desde el estado de juego en un juego
de cartas coleccionables~\cite{janusz_hearthstone_2017}. Esta es una
revisi\'on individual; el art\'iculo conjunto de la Etapa 1 la fusionar\'a
con las revisiones de los Ejes A y B escritas por los otros dos integrantes
del equipo, para un corpus combinado de al menos veintiun trabajos.

\subsection{Eje A: el problema y su dominio}
\label{subsec:axisa}
\emph{\textquestiondown C\'omo se ha abordado este problema en la literatura?
\textquestiondown Qu\'e enfoques dominan? \textquestiondown Qu\'e se considera resuelto y qu\'e sigue
abierto?}

\subsubsection{Predicci\'on del desenlace a partir de un \'unico estado}
La formulaci\'on dominante trata la predicci\'on del desenlace como
clasificaci\'on binaria supervisada sobre una instant\'anea de la
posici\'on, con la etiqueta tomada del resultado registrado de la partida
completa. El AAIA'17 Data Mining Challenge lo enuncia exactamente en esos
t\'erminos para \emph{Hearthstone}: los participantes recib\'ian la
descripci\'on de un estado de juego y entregaban la verosimilitud de que el
jugador uno ganara, puntuada mediante el \'area bajo la curva
ROC~\cite{janusz_hearthstone_2017}. El descriptor SC2EGSet aplica el mismo
encuadre a \emph{StarCraft II}, derivando indicadores econ\'omicos de las
repeticiones y clasificando el desenlace de la partida a partir de
ellos~\cite{bialecki_sc2egset_2023}. Ninguna de las dos formulaciones
requiere un simulador en el momento de la predicci\'on, y ninguna trata la
trayectoria como un problema de control; el estado es la entrada y el
desenlace registrado es la etiqueta. Esta es la formulaci\'on que se adopta
aqu\'i.

Lo que ambos trabajos establecen es que la formulaci\'on es viable en juegos
de esta familia y que la se\~nal es sustancial. Lo que no establecen es que
el nivel alcanzable se transfiera entre juegos. Los dos puntos de
operaci\'on reportados difieren m\'as entre s\'i que la diferencia entre
m\'etodos dentro de cualquiera de los dos estudios, y no est\'an medidos en
la misma escala: Janusz \textit{et al.} reportan AUC sobre estados
extra\'idos de partidas entre agentes que eligen sus acciones al azar,
mientras que Bia\l ecki \textit{et al.} reportan acierto sobre
caracter\'isticas promediadas a lo largo de toda la partida. Un estudio
sobre un tercer juego no puede heredar ninguna de las dos cifras como
expectativa.

\subsubsection{Capacidad predictiva a lo largo de la partida}
Un segundo tema, presente en ambos trabajos pero desarrollado por ninguno,
es que la capacidad predictiva no es un escalar. Janusz \textit{et al.}
grafican el AUC sobre turnos consecutivos y muestran que crece a lo largo
de la partida~\cite{janusz_hearthstone_2017}; Bia\l ecki \textit{et al.}
grafican el acierto sobre el tiempo de partida a intervalos de
aproximadamente siete segundos y observan que sus tres clasificadores
alcanzan una as\'intota alrededor de los 550 segundos, lo que interpretan
como el punto en que los indicadores econ\'omicos dejan de aportar
informaci\'on adicional y otros factores toman el
relevo~\cite{bialecki_sc2egset_2023}. Ambas curvas se\~nalan el mismo punto
metodol\'ogico: una \'unica cifra agregada promedia un r\'egimen en el que
el estado dice poco con otro en el que lo dice casi todo. Ninguno de los dos
trabajos reporta intervalos de confianza sobre esas curvas, y ninguno
reporta c\'omo cambia su forma cuando cambia la poblaci\'on de jugadores.

\subsubsection{Generalizaci\'on ante contenido cambiante}
El tercer tema es que el contenido de estos juegos no se queda quieto.
Bertram \textit{et al.} lo abordan directamente para \emph{Magic: The
Gathering}, donde se publican conjuntos de cartas nuevos con regularidad, y
eval\'uan reservando un conjunto de cartas completo. Sus resultados separan
dos capacidades que una \'unica puntuaci\'on dentro de la distribuci\'on
confunde. Entrenada sobre un solo conjunto, una codificaci\'on
\emph{one-hot} de carta alcanza un $67.80\%$ de acierto \emph{top-1}
sobre datos reservados de ese mismo conjunto y no puede producir
predicci\'on alguna para cartas que nunca ha visto, mientras que una
incrustaci\'on aleatoria de la misma dimensionalidad alcanza un
$67.87\%$ dentro de la distribuci\'on y un $23.79\%$ fuera de ella, que
es el acierto del azar para su tarea de
elecci\'on~\cite{bertram_cardrepresentations_2024}. Las representaciones
construidas a partir de atributos de carta, metadatos y codificaciones de
imagen alcanzan un $68.00\%$ dentro de la distribuci\'on y un
$42.87\%$ sobre conjuntos no vistos; escalado a un corpus de 2\,990
cartas y unos 75 millones de decisiones, el mismo modelo alcanza un
$55.44\%$ sobre el conjunto \textsc{bro}, enteramente
reservado~\cite{bertram_cardrepresentations_2024}. La lecci\'on trasciende
su tarea: una codificaci\'on del contenido basada en la identidad puede ser
indistinguible de una sem\'antica dentro de la distribuci\'on y no valer
nada fuera de ella.

\subsubsection{Qu\'e est\'a resuelto y qu\'e sigue abierto}
Resuelto: la predicci\'on del desenlace a partir de un estado intermedio es
un problema supervisado bien planteado y tratable en juegos de cartas
coleccionables y en juegos de estrategia en tiempo real; las
caracter\'isticas tabulares derivadas del estado transportan se\~nal
sustancial; y la capacidad predictiva var\'ia sistem\'aticamente a lo largo
de la partida. Abierto: si las probabilidades estimadas est\'an calibradas y
no meramente bien ordenadas; c\'omo se comporta el perfil por etapa bajo un
cambio declarado en la poblaci\'on de jugadores o de mazos; y si algo de
esto se sostiene en el \emph{Pok\'emon TCG}, para el que no existe estudio
publicado.

\subsection{Eje B: t\'ecnicas del track seleccionado}
\label{subsec:axisb}
\emph{\textquestiondown Cu\'al es el estado actual de los m\'etodos que pensamos
utilizar, cu\'ales son sus fortalezas y l\'imites, y bajo qu\'e supuestos
funcionan?}

\subsubsection{El rendimiento marginal de la complejidad del modelo}
La evidencia disponible aqu\'i sostiene que los modelos tabulares
cl\'asicos son el instrumento adecuado para esta tarea, y lo hace con
cifras y no con preferencias. En la competencia de \emph{Hearthstone} se
registraron 296 equipos y 188 enviaron soluci\'on; todos los equipos mejor
clasificados usaron redes neuronales y el ganador us\'o redes
convolucionales, y aun as\'i la puntuaci\'on ganadora de 0.8019 de AUC
supera al \emph{baseline} de los organizadores, 0.7846, una red
totalmente conectada con dos capas ocultas entrenada sobre la
representaci\'on tabular, por menos de un dos por ciento, y los
organizadores lo hacen notar expl\'icitamente~\cite{janusz_hearthstone_2017}.
La regresi\'on log\'istica aparece entre los enfoques que describen como
simples y no obstante eficaces. En SC2EGSet, tres clasificadores
convencionales con hiperpar\'ametros iniciales t\'ipicos quedan a menos de
0.016 de acierto unos de otros en la tarea de dos jugadores: m\'aquina de
vectores de soporte con n\'ucleo RBF en 0.9071, XGBoost en 0.8924 y
regresi\'on log\'istica en 0.8916~\cite{bialecki_sc2egset_2023}. Dos
estudios independientes sobre dos juegos distintos coinciden, por tanto, en
que la mayor parte de la se\~nal extra\'ible de una representaci\'on
tabular del estado de juego es alcanzable sin arquitecturas profundas, y en
que el residuo es lo bastante peque\~no como para que la interpretabilidad
y el costo dominen la elecci\'on.

Una tercera fuente afina el argumento en direcci\'on a los modelos m\'as
simples. Karbalaie \textit{et al.} compararon diez clasificadores y
encontraron que los modelos lineales fueron los m\'as honestos bajo
validaci\'on consciente del participante: la regresi\'on log\'istica y el
an\'alisis discriminante lineal difirieron a lo sumo en 0.03 de acierto
entre el protocolo contaminado y el correcto, y mostraron brechas entre
entrenamiento y prueba de aproximadamente 0.08, mientras que los
\emph{ensembles}, \emph{extra trees}, \emph{random forest},
\emph{gradient boosting}, mostraron brechas de al menos 0.34 y fueron
los modelos cuyas puntuaciones m\'as infl\'o el protocolo
defectuoso~\cite{karbalaie_participantaware_2026}. Los modelos de alta
capacidad no son solo m\'as dif\'iciles de interpretar: son los modelos a
los que un protocolo defectuoso favorece. Ese es un argumento para mantener
un modelo lineal regularizado en la comparaci\'on como punto de referencia,
y no solo por formalidad.

\subsubsection{Validaci\'on cuando las observaciones no son independientes}
El supuesto que estos m\'etodos comparten es que las muestras son
independientes, y es exactamente el supuesto que este corpus viola.
Karbalaie \textit{et al.} cuantifican la consecuencia sobre 623 ensayos de
72 participantes: la validaci\'on cruzada estratificada de diez pliegues da
un acierto de 0.91 para \emph{extra trees}, y el mismo modelo bajo
validaci\'on dejando un participante fuera da
0.66~\cite{karbalaie_participantaware_2026}. Su explicaci\'on es la que se
transfiere: en datos de medidas repetidas el tama\~no muestral efectivo se
aproxima al n\'umero de participantes, no al n\'umero de ensayos.
Sustituir participante por episodio e instant\'anea por ensayo enuncia la
restricci\'on que gobierna el presente trabajo. Reportan adem\'as que la
validaci\'on cruzada agrupada de $k$ pliegues coincide con el dise\~no
anidado y con dejar uno fuera dentro de 0.02 de acierto a un costo mucho
menor, lo que importa cuando la unidad de agrupaci\'on se cuenta en decenas
de miles. \textbf{Brecha de dominio, declarada expl\'icitamente:} se trata
de un estudio de biomec\'anica, y se cita por su argumento estad\'istico
sobre observaciones dependientes, no como resultado del dominio de juegos.

\subsubsection{La fuga de datos como objeto formal}
Kaufman \textit{et al.} aportan el marco que vuelve verificables, y no
meramente intuitivas, esas decisiones. Definen la fuga de datos (\emph{data
leakage}) como una propiedad intr\'inseca de la relaci\'on entre las
entradas observacionales de un modelo y la instancia objetivo: una
caracter\'istica es leg\'itima solo si es observable para el prop\'osito de
inferir ese objetivo en particular, evaluada instancia por instancia. Dos
consecuencias gobiernan nuestro dise\~no. La primera es el \emph{requisito
de la no m\'aquina del tiempo}: un modelo leg\'itimo solo puede usar
informaci\'on de un instante no posterior al del objetivo. La segunda es la
\emph{separaci\'on aprender--predecir}, la metodolog\'ia de evitaci\'on en
la que cada observaci\'on se etiqueta con la informaci\'on necesaria para
decidir su legitimidad y el conjunto de entrenamiento se construye despu\'es
de modo que solo entren entradas leg\'itimas~\cite{kaufman_leakage_2011}.
Advierten adem\'as que la fuga entra por las decisiones de dise\~no, 
dise\~no de caracter\'isticas, selecci\'on de modelo, elecci\'on de
hiperpar\'ametros informada por conocimiento que el modelo desplegado no
tendr\'ia, a la que llaman la forma m\'as peligrosa, porque un evaluador
no puede detectarla ni siquiera bus\-c\'an\-do\-la, y se\~nalan que eliminar
las fugas evidentes puede no sellar el resto.

Sasse \textit{et al.} actualizan ese marco convirti\'endolo en un
cat\'alogo de escenarios que los profesionales encuentran realmente,
organizado como fuga de prueba a entrenamiento, de prueba a prueba, de
caracter\'istica a objetivo, de objetivo, de conjunto de datos y de
variables de confusi\'on~\cite{sasse_leakage_2025}. Tres de sus entradas se
corresponden directamente con decisiones de este trabajo. Su tratamiento de
la fuga de prueba a entrenamiento enumera el no tener en cuenta muestras
dependientes como causa primaria y lo demuestra sobre pares de hermanos en
una cohorte de neuroimagen, donde permitir que familiares caigan a ambos
lados de la partici\'on eleva el acierto medido. Su caso de preprocesamiento
establece que el escalado, la imputaci\'on, la reducci\'on de
dimensionalidad y el sobremuestreo estimados sobre el conjunto completo
antes de particionar constituyen por s\'i mismos fuga de datos, raz\'on por
la cual cualquier correcci\'on de balance de clases en este trabajo debe
ajustarse dentro del pliegue de entrenamiento. Su categor\'ia de
caracter\'istica a objetivo cubre el caso que aqu\'i m\'as importa:
variables informativas y disponibles durante el aprendizaje pero ausentes en
el despliegue. Su prescripci\'on es una pregunta antes que una regla, 
comprobar si cada caracter\'istica estar\'ia disponible tras el despliegue, y es la pregunta que aplicamos campo por campo en la
Secci\'on~\ref{subsec:leakage}.

\subsubsection{Justificaci\'on t\'ecnica del Track A}
El Track A se elige por la fuerza de esa evidencia y no por descarte. El
corpus de repeticiones se reduce de forma natural a una tabla de ancho fijo
con campos num\'ericos, binarios y categ\'oricos, que es el r\'egimen en el
que los dos estudios del dominio de juegos citados arriba muestran una
paridad casi total entre modelos cl\'asicos y profundos. La pregunta
cient\'ifica es cu\'ales caracter\'isticas del estado transportan
informaci\'on sobre la victoria, lo que exige modelos cuyas atribuciones
puedan inspeccionarse. Y el riesgo dominante es de protocolo y no de
capacidad: la contaminaci\'on medida de una partici\'on a nivel de fila hace
que sea el dise\~no de la partici\'on, y no la arquitectura, la variable que
determina si un resultado tiene significado. Un modelo secuencial profundo
sobre la trayectoria, la alternativa natural del Track B, responder\'ia
una pregunta distinta y m\'as amplia, heredando todos los riesgos de fuga y
ninguna de las ventajas de interpretabilidad. El aprendizaje por refuerzo,
la lectura natural de un corpus de episodios, se descarta porque el objeto
de estudio es la funci\'on de valor de forma aislada, aprendida
\emph{offline} a partir de desenlaces registrados, y no una pol\'itica.

\subsection{Eje C: conjunto de datos y conjuntos an\'alogos}
\label{subsec:axisc}
\emph{\textquestiondown Qu\'e trabajos han utilizado el conjunto de datos, bajo qu\'e
protocolo, cu\'al es el mejor resultado publicado y qu\'e problemas
presenta?}

\subsubsection{Sin literatura previa sobre este corpus}
No localizamos ninguna publicaci\'on arbitrada que utilice los episodios del
PTCG AI Battle Challenge~\cite{ptcg_dataset_2026}. El corpus se public\'o en
2026 y est\'a ligado a una competencia todav\'ia en curso. De acuerdo con el
enunciado del curso, el eje se cubre por tanto con conjuntos de datos
an\'alogos del mismo dominio, y la sustituci\'on se argumenta en lugar de
darse por supuesta.

\subsubsection{Por qu\'e los an\'alogos son comparables, y d\'onde no lo son}
SC2EGSet es el an\'alogo estructural m\'as cercano: un corpus de
repeticiones competitivas de un juego de dos jugadores, procesado en
registros de estado de juego por partida, publicado junto con sus
herramientas de extracci\'on y acompa\~nado de experimentos de predicci\'on
de victoria sobre caracter\'isticas
tabulares~\cite{bialecki_sc2egset_2023}. Comparte la forma de la tuber\'ia,  repetici\'on cruda, estado analizado, caracter\'isticas tabulares,
etiqueta supervisada de desenlace, y el hecho de que los datos fueron
producidos por competencia y no por experimento. Difiere en g\'enero, en
tanto \emph{StarCraft II} es un juego de estrategia en tiempo real con
tiempo continuo y no un juego de cartas por turnos con recursos discretos,
y difiere en unidad de observaci\'on, puesto que sus modelos consumen
agregados por partida y no estados instant\'aneos. El conjunto de datos de
la competencia de \emph{Hearthstone} es el an\'alogo m\'as cercano en
g\'enero: un juego de cartas coleccionables con mazos, mano, tablero e
informaci\'on oculta, cuya tarea es precisamente la predicci\'on del ganador
a partir de un estado intermedio~\cite{janusz_hearthstone_2017}. Difiere en
la calidad de los agentes que lo generaron y en que sus estados fueron
publicados por los organizadores en formato tabular y JSON crudo en lugar de
ser analizados por quien modela. El corpus de \emph{drafting} de
\emph{Magic: The Gathering} es an\'alogo solo en el sentido estrecho que
importa para una de nuestras decisiones: es un juego de cartas
coleccionables cuyo contenido cambia con el tiempo y cuya evaluaci\'on
reserva contenido no visto~\cite{bertram_cardrepresentations_2024}. Su tarea
es predecir una elecci\'on humana de carta, no un desenlace de partida, y no
debe leerse como un resultado de predicci\'on de victoria.

\subsubsection{Problemas documentados de los corpus comparables}
Los problemas que estos autores declaran son los problemas que deber\'iamos
esperar. Bia\l ecki \textit{et al.} afirman que su corpus contiene
\'unicamente repeticiones de nivel de torneo y que, en consecuencia, ``no
existe posibilidad de realizar an\'alisis m\'as generales que tomen en
consideraci\'on a personas de distintos niveles de habilidad''; que los
cambios del juego a lo largo del tiempo hacen que algunos campos de datos
tomen valor cero, de modo que la representaci\'on del juego cambi\'o con el
tiempo; que la completitud dependi\'o de qu\'e organizadores de torneo
decidieron publicar; y que deliberadamente no filtraron repeticiones
terminadas de forma prematura, con el fin de preservar las distribuciones
originales~\cite{bialecki_sc2egset_2023}. Cada uno de esos puntos tiene su
contraparte en nuestro corpus: los episodios se seleccionan por puntuaci\'on
media de los agentes y se truncan a un tope diario de volumen; seis
versiones del motor coexisten en la muestra auditada; y los episodios
abandonados est\'an presentes y son identificables.

Janusz \textit{et al.} documentan dos problemas propios. Sus estados
proven\'ian de partidas entre agentes que eleg\'ian sus decisiones dentro del
juego al azar, lo que constituye una poblaci\'on espec\'ifica y no
representativa; y se detect\'o una fuga de datos involuntaria durante las
primeras semanas de la competencia, que oblig\'o a los organizadores a
mover 1\,250\,000 casos del conjunto de prueba al de
entrenamiento~\cite{janusz_hearthstone_2017}. Vale la pena enunciar lo
segundo con claridad: un incidente de fuga dentro de una competencia
organizada profesionalmente sobre exactamente esta tarea es la evidencia
m\'as fuerte disponible de que el riesgo es real y no hipot\'etico.

\subsubsection{Cambio de distribuci\'on y qu\'e da\~na primero}
Su respuesta parcial al problema de la poblaci\'on es tambi\'en el resultado
de generalizaci\'on m\'as \'util del que disponemos. Habiendo entrenado
sobre estados de juego aleatorio, construyeron un segundo corpus de 28\,604
partidas y 814\,407 estados generados por agentes fuertes de b\'usqueda en
\'arbol Monte Carlo que realizaban 15\,000 simulaciones por acci\'on, y
observaron que el AUC total ca\'ia solo de aproximadamente 0.79 a
aproximadamente 0.75 al evaluar sobre esa poblaci\'on
desplazada~\cite{janusz_hearthstone_2017}. Un cambio en la calidad de los
agentes que generan los estados cuesta, por tanto, algo, pero mucho menos de
lo que podr\'ia temerse.

Lo que ese experimento no mide es la calibraci\'on, y la evidencia
sistem\'atica indica que esa es la magnitud en mayor riesgo. Guo
\textit{et al.} examinaron 4\,457 referencias y retuvieron quince estudios
que cuantificaban el cambio temporal del conjunto de datos e implementaban
alguna mitigaci\'on. Los once estudios que evaluaron calibraci\'on
reportaron deterioro de la calibraci\'on, mientras que solo tres de los doce
que evaluaron discriminaci\'on reportaron deterioro de la discriminaci\'on;
los tres que lo hicieron eran estudios de un solo centro, mientras que los
nueve que no reportaron deterioro eran
multic\'entricos~\cite{guo_temporalshift_2021}. Su inventario de
mitigaciones est\'a dominado por el reajuste del modelo ($n=12$), la
calibraci\'on de probabilidades ($n=7$), la actualizaci\'on del modelo
($n=6$) y la selecci\'on de modelo ($n=6$), frente a solo dos enfoques a
nivel de caracter\'isticas, y su s\'intesis es que todas esas estrategias
preservaron la calibraci\'on pero ninguna preserv\'o uniformemente la
discriminaci\'on. Aportan adem\'as el vocabulario que este trabajo
necesita, distinguiendo el cambio de covariables
($P_{\text{ent}}(x)\neq P_{\text{prueba}}(x)$ con
$P_{\text{ent}}(y|x)=P_{\text{prueba}}(y|x)$) del cambio de concepto
($P_{\text{ent}}(y|x)\neq P_{\text{prueba}}(y|x)$ con
$P_{\text{ent}}(x)=P_{\text{prueba}}(x)$), y se\~nalando que ambos pueden
coexistir. \textbf{Brecha de dominio, declarada expl\'icitamente:} se trata
de predicci\'on cl\'inica, citada por la asimetr\'ia que mide y por la
terminolog\'ia que estandariza, no como hallazgo del dominio de juegos. Sus
propias limitaciones declaradas, un barrido incompleto de actas de
congresos, modelos desplegados no publicados y la ausencia de tratamiento
del cambio geogr\'afico, no afectan a ninguno de los dos usos.

% TABLA COMPARATIVA OBLIGATORIA (integrada; antes en tab/comparative_table)
\subsection{An\'alisis comparativo}
\label{subsec:comparative}

% ============================================================
% TABLA COMPARATIVA — PARTE 1 DE 2
% ============================================================

\begin{table*}[!tp]
\centering
\caption{Resumen comparativo de los siete trabajos revisados. Las filas se
ordenan por proximidad a nuestra tarea. Los valores est\'an transcritos de
la fuente; ``n/a'' marca un campo que el trabajo no posee, no un campo que
no hayamos podido encontrar.}
\label{tab:comparative}

\footnotesize
\setlength{\tabcolsep}{3.5pt}

\begin{tabularx}{\textwidth}{@{}l c P{2.05cm} P{2.35cm} P{1.5cm} P{2.35cm} Y@{}}
\toprule
\textbf{Referencia} & \textbf{A\~no} & \textbf{T\'ecnica} &
\textbf{Conjunto de datos / entorno} & \textbf{M\'etrica principal} &
\textbf{Resultado reportado} & \textbf{Limitaci\'on identificada}\\
\midrule

\cite{janusz_hearthstone_2017} & 2017 &
Competencia: \emph{ensembles} de CNN y ANN, XGBoost, regresi\'on
log\'istica; el \emph{baseline} de los organizadores es una red totalmente
conectada de 2 capas ocultas sobre datos tabulares &
Estados de juego de \emph{Hearthstone}; 3.25\,M de estados de entrenamiento
provenientes de $\approx$65\,000 partidas, 750\,000 estados de prueba
provenientes de $>$180\,000 partidas; 9 mazos en entrenamiento frente a 27
en prueba &
ROC-AUC &
0.7846 el \emph{baseline}; 0.8019 la mejor soluci\'on; brecha $<2\%$.
Transferencia a estados de agentes MCTS: $\approx$0.79 $\rightarrow$
$\approx$0.75 &
Estados generados por agentes que eligen acciones al azar, una poblaci\'on
no representativa; no se reporta calibraci\'on; se detect\'o una fuga
involuntaria a mitad de competencia que oblig\'o a pasar 1.25\,M de casos de
prueba al conjunto de entrenamiento\\
\addlinespace

\cite{bialecki_sc2egset_2023} & 2023 &
Regresi\'on log\'istica, SVM (RBF) y XGBoost con hiperpar\'ametros
pr\'oximos a los predeterminados; validaci\'on cruzada de 5 pliegues &
SC2EGSet: 55 paquetes de repeticiones de torneos de \emph{StarCraft II};
11\,115 partidas \'unicas reducidas a 22\,230 muestras promediadas; 17
caracter\'isticas econ\'omicas seleccionadas &
Acierto &
0.9071 con dos jugadores (SVM-RBF), 0.8924 XGBoost, 0.8916 regresi\'on
log\'istica; 0.8488 con un jugador &
Solo jugadores de nivel de torneo, por lo que no admite an\'alisis entre
niveles de habilidad; los autores llaman a sus experimentos ``m\'as
ilustrativos que investigativos''; campos puestos a cero conforme el juego
cambi\'o; no se describe agrupaci\'on por partida para los pliegues\\
\addlinespace

\cite{bertram_cardrepresentations_2024} & 2024 &
Redes siamesas bajo ordenamiento contextual de preferencias, sobre
representaciones generalizadas de carta (atributos, metadatos, im\'agenes
autocodificadas) &
Elecciones de \emph{draft} de \emph{Magic: The Gathering};
$\approx$100\,M de elecciones en 14 conjuntos de cartas; un conjunto
completo reservado &
Acierto \emph{top-1} &
67.20\% sobre conjuntos vistos frente a 55.44\% sobre el conjunto no
visto; la codificaci\'on \emph{one-hot} queda indefinida sobre cartas no
vistas; la incrustaci\'on aleatoria da 23.79\% fuera de la distribuci\'on &
La tarea es predecir una elecci\'on humana de carta y no un desenlace de
partida, por lo que solo es an\'aloga en cuanto a generalizaci\'on de
contenido; las elecciones humanas no son ni correctas ni consistentes, lo
que acota el acierto alcanzable; el ajuste fino sobre un conjunto degrada
los dem\'as\\
\addlinespace

\cite{karbalaie_participantaware_2026} & 2026 &
Diez clasificadores comparados bajo cuatro dise\~nos de validaci\'on:
estratificada de 10 pliegues, dejando un participante fuera, agrupada de 3
pliegues y anidada (dejando uno fuera $+$ bucle interno agrupado) &
623 ensayos de salto de 72 participantes tras reconstrucci\'on de ligamento
cruzado anterior; 624 caracter\'isticas extra\'idas reducidas a un
\emph{top-10} fijo &
Acierto, brecha entrenamiento--prueba, estabilidad de rangos &
\emph{Extra trees} 0.91 bajo estratificada de 10 pliegues frente a 0.66
dejando un participante fuera; los modelos lineales difieren $\leq$0.03; la
anidada coincide con dejar uno fuera dentro de 0.02 &
Dominio de biomec\'anica, citado por su argumento estad\'istico y no como
resultado del dominio de juegos; una \'unica tarea binaria sobre 72
participantes; calibraci\'on y atribuci\'on de caracter\'isticas
expl\'icitamente fuera de alcance\\

\bottomrule
\end{tabularx}
\end{table*}


% ============================================================
% TABLA COMPARATIVA — PARTE 2 DE 2
% ============================================================

\begin{table*}[!tp]

% Mantiene el mismo número de tabla que la parte anterior.
\addtocounter{table}{-1}

\centering
\caption{Resumen comparativo de los siete trabajos revisados
(continuaci\'on).}
\label{tab:comparative-cont}

\footnotesize
\setlength{\tabcolsep}{3.5pt}

\begin{tabularx}{\textwidth}{@{}l c P{2.05cm} P{2.35cm} P{1.5cm} P{2.35cm} Y@{}}
\toprule
\textbf{Referencia} & \textbf{A\~no} & \textbf{T\'ecnica} &
\textbf{Conjunto de datos / entorno} & \textbf{M\'etrica principal} &
\textbf{Resultado reportado} & \textbf{Limitaci\'on identificada}\\
\midrule

\cite{guo_temporalshift_2021} & 2021 &
Revisi\'on sistem\'atica (PRISMA) de estrategias de mitigaci\'on: reajuste
del modelo ($n$=12), calibraci\'on de probabilidades ($n$=7),
actualizaci\'on del modelo ($n$=6), selecci\'on de modelo ($n$=6) &
15 estudios de predicci\'on cl\'inica retenidos de 4\,457 referencias
examinadas; regresi\'on log\'istica usada en 13 de 15 &
Cambio en calibraci\'on y en discriminaci\'on &
Deterioro de calibraci\'on en 11 de los 11 estudios que la midieron;
deterioro de discriminaci\'on en solo 3 de 12 &
Dominio cl\'inico, citado por la asimetr\'ia medida y por la terminolog\'ia
del cambio de distribuci\'on; no se realiz\'o metaan\'alisis; cobertura
incompleta de actas de congresos; no se examin\'o el cambio geogr\'afico\\
\addlinespace

\cite{sasse_leakage_2025} & 2025 &
Cat\'alogo de escenarios de fuga: de prueba a entrenamiento, de prueba a
prueba, de caracter\'istica a objetivo, de objetivo, decaimiento del
conjunto de datos y de variables de confusi\'on, con demostraciones
ilustrativas &
n/a --- art\'iculo de revisi\'on; las demostraciones usan una cohorte de
neuroimagen, sin corpus de referencia propio &
n/a --- sin m\'etrica principal \'unica &
n/a --- los efectos se muestran gr\'aficamente por escenario en lugar de
tabularse como una cifra \'unica &
Es un cat\'alogo y no una comparaci\'on cuantitativa, por lo que establece
qu\'e fallos existen pero no su magnitud relativa; las demostraciones se
limitan a un solo tipo de dato\\
\addlinespace

\cite{kaufman_leakage_2011} & 2011 &
Definici\'on formal de la fuga mediante la legitimidad de las entradas
observacionales; evitaci\'on por etiquetado de legitimidad y separaci\'on
aprender--predecir; detecci\'on por an\'alisis exploratorio, an\'alisis
posterior y prueba temprana en campo &
n/a --- trabajo metodol\'ogico; ilustrado con competencias publicadas de
miner\'ia de datos &
n/a &
n/a --- no se reporta experimento alguno &
Sin evaluaci\'on emp\'irica de la metodolog\'ia propuesta; los autores
afirman que los intentos de reparar la fuga a posteriori suelen ser
incompletos y pueden empeorar la situaci\'on\\

\bottomrule
\end{tabularx}
\end{table*}

De la Tabla~\ref{tab:comparative} emergen tres patrones. El primero es una
divisi\'on de m\'etricas que bloquea la comparaci\'on directa: los dos
estudios sobre estados de juego reportan escalares distintos, AUC en un caso
y acierto en el otro, y ninguno reporta una regla de puntuaci\'on propia ni
una curva de fiabilidad. Cualquier afirmaci\'on de que uno de estos sistemas
est\'a mejor calibrado que otro carecer\'ia de respaldo, porque ninguno lo
midi\'o. El segundo es que los estudios que miden generalizaci\'on bajo
cambio la miden contra cosas diferentes, calidad de los
agentes~\cite{janusz_hearthstone_2017}, conjuntos de cartas no
vistos~\cite{bertram_cardrepresentations_2024}, identidad del
participante~\cite{karbalaie_participantaware_2026}, tiempo
calendario~\cite{guo_temporalshift_2021}, y que las degradaciones
reportadas son en todos los casos lo bastante grandes como para importar:
0.79 $\rightarrow$ 0.75 de AUC, 67.20\% $\rightarrow$ 55.44\% de
acierto, 0.91 $\rightarrow$ 0.66 de acierto. Una \'unica cifra dentro de la
distribuci\'on es la magnitud menos informativa que reporta cualquiera de
estos trabajos.

El tercer patr\'on es una limitaci\'on recurrente antes que una tendencia.
En todas las entradas emp\'iricas, la unidad de observaci\'on es menor que
la unidad que realmente es independiente, y los trabajos difieren en si
act\'uan en consecuencia. Janusz \textit{et al.} separan entrenamiento y
prueba a nivel de partidas completas, de modo que su protocolo est\'a
agrupado aunque no lo formulen as\'i. Karbalaie \textit{et al.} convierten
la agrupaci\'on en el objeto de estudio. Bia\l ecki \textit{et al.} reportan
validaci\'on cruzada de cinco pliegues sobre 22\,230 muestras promediadas
derivadas de 11\,115 partidas \'unicas, dos filas por partida, una por
jugador, sin describir restricci\'on alguna de agrupaci\'on a nivel de
partida. \textbf{Esta es una inferencia nuestra, no una afirmaci\'on de los
autores:} si esos pliegues se extrajeron a nivel de fila, ambas filas de una
partida podr\'ian caer a lados opuestos de una partici\'on, y sus aciertos
reportados ser\'ian optimistas en una cantidad desconocida. Lo se\~nalamos
como salvedad de comparabilidad y no como error, puesto que el art\'iculo no
declara el particionador que utiliz\'o.

Una cuarta observaci\'on concierne al peso de la evidencia. Dos de los siete
trabajos son referencias metodol\'ogicas sin conjunto de datos ni m\'etrica
propios~\cite{kaufman_leakage_2011,sasse_leakage_2025}; sus celdas se marcan
como no aplicables en lugar de rellenarse con un valor plausible. Forzar
n\'umeros en esas filas tergiversar\'ia lo que esos trabajos son.

% ===========================================================================
% 6.2.4  BASELINE DE LITERATURA
% ===========================================================================
\subsection{\emph{Baseline} de literatura}
\label{subsec:litbaseline}

No existe resultado publicado sobre el corpus objetivo, de modo que el
\emph{baseline} se toma del trabajo an\'alogo cuya tarea, unidad de
observaci\'on y m\'etrica son las m\'as cercanas a las nuestras.

\begin{quote}
\textbf{Objetivo de referencia:} ROC-AUC en la banda $0.78$--$0.80$,
proveniente del AAIA'17 Data Mining Challenge sobre \emph{Hearthstone}.
Tarea: predecir el ganador a partir de un \'unico estado de juego. Unidad de
observaci\'on: un estado de juego. M\'etrica: \'area bajo la curva ROC.
\emph{Baseline} oficial de los organizadores: $0.7846$, obtenido por una red
totalmente conectada con dos capas ocultas sobre la representaci\'on
tabular, que qued\'o en el puesto $94$. Mejor soluci\'on enviada: $0.8019$,
un \emph{ensemble} construido sobre redes convolucionales. Protocolo:
3\,250\,000 estados de entrenamiento provenientes de aproximadamente
65\,000 partidas frente a 750\,000 estados de prueba provenientes de m\'as
de 180\,000 partidas distintas, con 9 mazos en entrenamiento y 27 en prueba,
de modo que el conjunto de prueba contiene cartas nunca vistas en
entrenamiento; puntuaciones preliminares sobre un $5\%$ aleatorio fijo del
conjunto de prueba y puntuaciones finales sobre el
resto~\cite{janusz_hearthstone_2017}.
\end{quote}

\noindent\textbf{Qu\'e pretende hacer este trabajo con esa cifra.} El
objetivo operativo es que el modelo cl\'asico m\'as fuerte alcance esa banda
de $0.78$--$0.80$ de AUC bajo nuestro propio protocolo primario en la
ventana tard\'ia de la partida, y establecer en qu\'e punto de la partida se
alcanza por primera vez esa banda. No afirmamos que superaremos $0.8019$,
porque las dos cifras no est\'an medidas sobre el mismo objeto.

\noindent\textbf{Diferencias de protocolo que impiden una comparaci\'on
justa.}
\begin{itemize}
  \item \textbf{Juego distinto.} \emph{Hearthstone} y el \emph{Pok\'emon
        TCG} difieren en sistemas de recursos, condiciones de victoria y
        estructura de tablero, de modo que el contenido informativo de un
        estado no es la misma magnitud.
  \item \textbf{Poblaci\'on de agentes distinta.} Sus estados provienen de
        agentes que act\'uan al azar; los nuestros, de agentes que compiten
        y son seleccionados por puntuaci\'on. Su propio experimento de
        transferencia muestra que esto importa, y cuesta aproximadamente
        $0.04$ de AUC en la direcci\'on de agentes aleatorios a agentes
        fuertes~\cite{janusz_hearthstone_2017}.
  \item \textbf{Partici\'on distinta.} Su separaci\'on es por partida y por
        mazo entre dos conjuntos fijos; la nuestra es una partici\'on
        agrupada por episodio, con protocolos adicionales temporal y de
        reserva de mazos.
  \item \textbf{Prevalencia de clases distinta.} Su conjunto de
        entrenamiento tiene $50.46\%$ de victorias y el de prueba
        $57.27\%$, de modo que las dos mitades de su propio banco de
        pruebas no provienen de la misma
        distribuci\'on~\cite{janusz_hearthstone_2017}.
  \item \textbf{Descripci\'on del estado distinta.} Su representaci\'on
        tabular fue suministrada por los organizadores, y las mejores
        participaciones reconstruyeron caracter\'isticas a partir del JSON
        crudo; nuestro conjunto de caracter\'isticas lo construimos nosotros
        y est\'a restringido por una lista negra de fuga.
\end{itemize}

\noindent Se dispone de un segundo anclaje, m\'as d\'ebil. SC2EGSet reporta
$0.9071$ de acierto en predicci\'on de victoria con dos jugadores y
$0.8488$ con un solo jugador~\cite{bialecki_sc2egset_2023}. Se registra
aqu\'i por completitud y no se adopta como objetivo: la m\'etrica es acierto
y no AUC, la unidad es una partida completa y no un instante, las
caracter\'isticas son promedios sobre la partida, y la cuesti\'on de la
agrupaci\'on planteada en la Secci\'on~\ref{subsec:comparative} queda sin
resolver.

\subsection{Brechas identificadas}
\label{subsec:gaps}

\textbf{Brechas encontradas en la literatura.}
\begin{enumerate}
  \item \textbf{Ning\'un trabajo arbitrado utiliza este corpus.} Los
        episodios del PTCG AI Battle Challenge no tienen descripci\'on
        publicada, ni diccionario de datos documentado, ni \emph{baseline}
        reportado.
  \item \textbf{La tarea no se ha estudiado en el \emph{Pok\'emon TCG}.} La
        predicci\'on de victoria desde el estado de juego est\'a establecida
        en \emph{Hearthstone}~\cite{janusz_hearthstone_2017} y en
        \emph{StarCraft II}~\cite{bialecki_sc2egset_2023}, pero no en un
        juego cuya estructura de recursos se apoya en cartas de premio,
        asignaci\'on de energ\'ia, banca y evoluci\'on.
  \item \textbf{Se reporta discriminaci\'on sin calibraci\'on.} Ninguno de
        los dos estudios sobre estados de juego reporta una regla de
        puntuaci\'on propia ni una curva de
        fiabilidad~\cite{janusz_hearthstone_2017,bialecki_sc2egset_2023}, y
        Karbalaie \textit{et al.} colocan expl\'icitamente la calibraci\'on
        fuera de su alcance y la recomiendan como trabajo
        futuro~\cite{karbalaie_participantaware_2026}. Guo \textit{et al.}
        muestran que esta es la magnitud que se degrada
        primero~\cite{guo_temporalshift_2021}. Un modelo pensado como
        evaluador de posiciones est\'a siendo validado, por tanto, sobre la
        propiedad equivocada.
  \item \textbf{La evaluaci\'on agrupada se aplica de forma inconsistente en
        corpus de estados de juego.} El costo medido de ignorar la
        dependencia solo est\'a disponible desde otros
        campos~\cite{karbalaie_participantaware_2026,sasse_leakage_2025};
        dentro de los conjuntos de datos de juegos la restricci\'on unas
        veces se satisface sin ser
        nombrada~\cite{janusz_hearthstone_2017} y otras queda sin
        describir~\cite{bialecki_sc2egset_2023}.
  \item \textbf{La generalizaci\'on rara vez se eval\'ua bajo un cambio
        declarado.} Cuando se hace, la degradaci\'on es
        sustancial~\cite{janusz_hearthstone_2017,%
        bertram_cardrepresentations_2024}, y la evidencia sistem\'atica
        encuentra solo quince estudios que a la vez cuantifiquen el cambio
        temporal y act\'uen sobre \'el~\cite{guo_temporalshift_2021}.
\end{enumerate}

\noindent\textbf{Brechas que este trabajo aborda.} Las brechas 2 a 5. La 2,
planteando la tarea en el \emph{Pok\'emon TCG} por primera vez y reportando
el perfil por turno de la capacidad predictiva. La 3, reportando
\emph{log loss}, \emph{Brier score} y curvas de fiabilidad junto al ROC-AUC,
y poniendo a prueba la hip\'otesis H3 de forma expl\'icita. La 4, agrupando
por \code{episode\_id} toda partici\'on, todo intervalo de confianza y todo
pliegue de validaci\'on cruzada, y reportando la contaminaci\'on medida que
producir\'ia la alternativa ingenua. La 5, evaluando bajo tres cambios
declarados, episodios no vistos, mazos no vistos y periodos temporales
posteriores, y reportando cada uno por separado en lugar de promediarlos.

\noindent\textbf{Brechas fuera de alcance.} La brecha 1 se documenta, no se
cierra: un proyecto de un semestre no puede establecer una referencia
arbitrada para el corpus, y este trabajo no pretende serlo. La
comparabilidad estricta con el banco de pruebas de \emph{Hearthstone}
tambi\'en queda fuera de alcance, por las razones enumeradas en la
Secci\'on~\ref{subsec:litbaseline}; las diferencias son estructurales y no
pueden eliminarse reejecutando ninguno de los dos experimentos. M\'as
all\'a de estas, no intentamos construir un agente competitivo, aprender una
pol\'itica, separar la deriva del metajuego del cambio de versi\'on del
motor, la muestra auditada confunde ambas causas, ni generalizar
ning\'un hallazgo a jugadores humanos.

% ===========================================================================
% 6.2.5  METODOLOGIA PROPUESTA
% ===========================================================================
\section{Metodolog\'ia Propuesta}
\label{sec:methodology}

Esta secci\'on describe lo que proponemos hacer. No se ha entrenado
ning\'un modelo y ning\'un resultado aqu\'i reportado es emp\'irico; las
\'unicas magnitudes medidas que aparecen son propiedades del corpus,
establecidas por la auditor\'ia que se describe en la
Secci\'on~\ref{sec:trackanalysis}.

\subsection{Formalizaci\'on del problema}
\label{subsec:formalization}

Sea $\mathcal{E}$ el conjunto de episodios, cada uno una partida completa
entre dos agentes. Un episodio $e \in \mathcal{E}$ produce una secuencia
ordenada de estados; de ella retenemos la subsecuencia de estados
\emph{frescos} del asiento que act\'ua, indexada por $t$. Cada estado
retenido produce una observaci\'on
\[
  (x_i,\; y_i,\; g_i,\; \tau_i,\; d_i),
\]
donde $x_i \in \mathcal{X} \subseteq \mathbb{R}^p$ es el vector de
caracter\'isticas construido a partir de la observaci\'on disponible para el
asiento activo en ese instante, $y_i \in \{0,1\}$ indica si el jugador desde
cuya perspectiva se construy\'o la fila gan\'o el episodio, $g_i \in
\mathcal{E}$ es el identificador del episodio, $\tau_i \in \mathbb{N}$ es el
\'indice de turno y $d_i$ la fecha de calendario. Los tres \'ultimos campos
se arrastran para agrupaci\'on, estratificaci\'on y an\'alisis; por
construcci\'on, nunca son componentes de $x_i$.

La tarea es clasificaci\'on binaria probabil\'istica: aprender
$f: \mathcal{X} \rightarrow [0,1]$ tal que $f(x) \approx P(y = 1 \mid x)$.
Puesto que la salida es una probabilidad y no una etiqueta, el objetivo de
entrenamiento es una regla de puntuaci\'on propia, la
log-verosimilitud negativa regularizada:
\begin{equation}
\label{eq:objective}
\begin{split}
\hat{f} = \arg\min_{f \in \mathcal{F}} \Big\{
  &-\frac{1}{N}\sum_{i=1}^{N}\big[\, y_i \log f(x_i)\\
  &+ (1-y_i)\log\big(1-f(x_i)\big) \,\big]
  + \lambda\,\Omega(f) \Big\},
\end{split}
\end{equation}
con $\Omega$ la penalizaci\'on de complejidad de la familia de modelos
$\mathcal{F}$ y $\lambda$ seleccionado sobre datos de validaci\'on.
Minimizar una regla de puntuaci\'on propia en lugar de una p\'erdida de
clasificaci\'on es una elecci\'on deliberada: es el objetivo cuyo \'optimo
es la probabilidad condicional verdadera, y la calibraci\'on es una de las
propiedades que este trabajo se propone medir.

Dos restricciones completan la formulaci\'on. La primera es la condici\'on
de legitimidad de Kaufman \textit{et al.}: todo componente de $x_i$ debe ser
observable para el prop\'osito de inferir $y_i$, evaluado instancia por
instancia, y en particular no puede depender de informaci\'on generada
despu\'es del turno $\tau_i$~\cite{kaufman_leakage_2011}. La segunda es una
restricci\'on de partici\'on. Para cualesquiera dos partes $A$ y $B$ de una
divisi\'on,
\begin{equation}
\label{eq:grouping}
  \{\, g_i : i \in A \,\} \;\cap\; \{\, g_i : i \in B \,\} \;=\; \emptyset ,
\end{equation}
de modo que ning\'un episodio aporte filas a m\'as de una parte. La
restricci\'on \eqref{eq:grouping} es el enunciado formal del argumento del
tama\~no muestral efectivo de Karbalaie \textit{et al.}: la unidad
independiente es el episodio, no la
fila~\cite{karbalaie_participantaware_2026}.

\subsection{Tuber\'ia propuesta}
\label{subsec:pipeline}

\begin{figure}[!t]
\centering
\footnotesize
\setlength{\fboxsep}{6pt}
\fbox{\begin{minipage}{0.90\columnwidth}
\ttfamily\raggedright
repetici\'on JSON\\[2pt]
$\vert$\, validaci\'on de esquema (firma de nivel superior,
\code{schema\_version})\\
$\vert$\, selecci\'on de estados frescos del asiento activo\\
$\vert$\, proyecci\'on a la perspectiva del jugador activo (propio / rival)\\
$\vert$\, filtro de lista negra $\rightarrow$ falla la construcci\'on si un
campo prohibido llega al vector de caracter\'isticas\\
$\vert$\, ingenier\'ia de caracter\'isticas (familias M0--M4)\\[3pt]
$\Downarrow$\\[3pt]
tabla: una fila por instant\'anea, m\'as \code{episode\_id}, \code{date} y
\code{turn} (solo agrupaci\'on)\\[2pt]
$\vert$\, partici\'on agrupada por \code{episode\_id} $\rightarrow$
entrenamiento / validaci\'on / prueba\\
$\vert$\, ajuste del preprocesamiento SOLO sobre el pliegue de
entrenamiento\\[3pt]
$\Downarrow$\\[3pt]
escalera de modelos B0--B4 $\rightarrow$ m\'etricas por ventana de turno,
m\'as intervalos de confianza bootstrap sobre episodios
\end{minipage}}
\caption{Tuber\'ia propuesta, de la repetici\'on cruda al modelo evaluado.
El filtro de lista negra y la partici\'on agrupada se imponen mediante
pruebas automatizadas, no por convenci\'on.}
\label{fig:pipeline}
\end{figure}

La Fig.~\ref{fig:pipeline} resume el flujo previsto. Dos decisiones de
dise\~no merecen justificaci\'on. Primera, la perspectiva de cada fila se
normaliza a \emph{propio} frente a \emph{rival} usando el \'indice de
asiento, de modo que el modelo aprenda relaciones entre magnitudes y no la
identidad de un asiento; la auditor\'ia encontr\'o que asiento y orden de
turno est\'an casi perfectamente confundidos en los d\'ias tard\'ios, lo que
de otro modo convertir\'ia la identidad del asiento en un atajo atractivo.
Segunda, los par\'ametros de preprocesamiento, escalado, codificaci\'on
de variables categ\'oricas, cualquier imputaci\'on, cualquier ponderaci\'on
de clases, se estiman dentro del pliegue de entrenamiento y despu\'es se
aplican a los datos reservados. Sasse \textit{et al.} clasifican la
alternativa, estimarlos sobre el conjunto completo antes de particionar,
como fuga de prueba a entrenamiento, aunque el modelo en s\'i se ajuste
despu\'es de la partici\'on~\cite{sasse_leakage_2025}.

\subsection{Protocolo experimental}
\label{subsec:protocol}

\subsubsection{Partici\'on de los datos}
El protocolo primario es una partici\'on $70/15/15$ \emph{agrupada por
episodio}, que satisface \eqref{eq:grouping}. La alternativa, una
partici\'on aleatoria sobre filas, se rechaza sobre bases medidas y no de
principio: en la muestra auditada, una partici\'on aleatoria $70/30$ de
instant\'aneas colocar\'ia el $99.9945\%$ de los episodios a ambos lados
simult\'aneamente. Se proponen dos protocolos secundarios. Un protocolo
temporal entrena sobre los d\'ias auditados m\'as tempranos, valida sobre
los d\'ias intermedios y prueba sobre los m\'as tard\'ios, implementando la
separaci\'on aprender--predecir que Kaufman \textit{et al.} describen como
predicci\'on sobre el futuro~\cite{kaufman_leakage_2011}. Un protocolo de
reserva de mazos aparta mazos completos o pares de mazos para el conjunto de
prueba, que es el dise\~no que Bertram \textit{et al.} emplean para separar
la memorizaci\'on de contenido conocido de la comprensi\'on del
contenido~\cite{bertram_cardrepresentations_2024}. Los tres responden
preguntas distintas y se reportan por separado, nunca promediados.

\subsubsection{Estrategia de validaci\'on}
Los hiperpar\'ametros se seleccionan con validaci\'on cruzada agrupada de
$k$ pliegues sobre la partici\'on de entrenamiento, tomando el episodio como
clave de agrupaci\'on; la partici\'on de prueba se toca una sola vez, al
final. Karbalaie \textit{et al.} reportan que la validaci\'on cruzada
agrupada de $k$ pliegues coincide con el dise\~no de dejar un grupo fuera,
mucho m\'as costoso, dentro de 0.02 de acierto y a un costo un orden de
magnitud menor, lo que resulta decisivo cuando la unidad de agrupaci\'on se
cuenta en decenas de miles~\cite{karbalaie_participantaware_2026}. Se
prefiere un dise\~no plenamente anidado all\'a donde el presupuesto
computacional lo permita, puesto que tanto Karbalaie \textit{et al.} como
Sasse \textit{et al.} identifican el sesgo de selecci\'on como una fuga
residual cuando los mismos pliegues sirven para ajustar y para
evaluar~\cite{karbalaie_participantaware_2026,sasse_leakage_2025}.

\subsubsection{M\'etricas}
Cada m\'etrica se incluye porque responde una pregunta que las dem\'as no
pueden responder; la Tabla~\ref{tab:metrics} indica cu\'al.

\begin{table}[!t]
\centering
\caption{M\'etricas de evaluaci\'on y su justificaci\'on.}
\label{tab:metrics}
\footnotesize
\begin{tabularx}{\columnwidth}{@{}P{1.9cm} Y@{}}
\toprule
\textbf{M\'etrica} & \textbf{Pregunta que responde, y por qu\'e no basta por s\'i sola}\\
\midrule
ROC-AUC \newline \textit{(principal)} &
\textquestiondown Ordena el modelo correctamente dos posiciones? Es independiente del
umbral y robusta ante desbalance moderado, y es la m\'etrica de la \'unica
tarea publicada comparable~\cite{janusz_hearthstone_2017}. No dice nada
sobre si la probabilidad es correcta.\\
\emph{Log loss} &
\textquestiondown Es exacta la probabilidad en s\'i misma? Es el objetivo de
entrenamiento \eqref{eq:objective}, de modo que reportarla muestra
adem\'as si la optimizaci\'on tuvo \'exito.\\
\emph{Brier score} &
\textquestiondown Cu\'al es el error cuadr\'atico de la probabilidad, y c\'omo se
descompone en calibraci\'on y refinamiento?\\
Curva de fiabilidad y error de calibraci\'on &
Entre las posiciones a las que el modelo asigna un $70\%$,
\textquestiondown gana efectivamente el $70\%$? Es la propiedad que Guo \textit{et al.}
encuentran que se degrada primero bajo cambio
temporal~\cite{guo_temporalshift_2021} y que ninguno de los dos estudios
comparables sobre estados de juego
reporta~\cite{janusz_hearthstone_2017,bialecki_sc2egset_2023}.\\
Acierto, precisi\'on, exhaustividad y F1 &
Se reportan solo una vez fijado y justificado un umbral, para comparabilidad
con trabajos que reportan acierto~\cite{bialecki_sc2egset_2023}.
Insuficientes por s\'i solas: un modelo puede acertar la clase y estar
sistem\'aticamente sobreconfiado.\\
\bottomrule
\end{tabularx}
\end{table}

ROC-AUC es la m\'etrica principal por tres razones: el objetivo est\'a
cercano al equilibrio a nivel de episodio, de modo que el acierto ser\'ia
poco informativo antes que enga\~noso; la m\'etrica no exige comprometerse
con un umbral, y en esta etapa no hay definido ning\'un costo de decisi\'on;
y es la \'unica m\'etrica sobre la que existe una cifra publicada
an\'aloga. Si la distribuci\'on de clases de la tabla final resultara
marcadamente desbalanceada, una magnitud a\'un no medida, puesto que
depende de cu\'antas instant\'aneas aporta cada bando, se a\~nadir\'a
PR-AUC, con la clase minoritaria declarada expl\'icitamente. Cada m\'etrica
se reporta de forma global y por ventana de turno, y todo intervalo de
confianza se obtiene por remuestreo \emph{bootstrap} \emph{sobre
episodios}, nunca sobre filas, por la misma raz\'on que gobierna
\eqref{eq:grouping}.

\subsubsection{Prevenci\'on de la fuga de datos}
Los vectores de fuga y sus contramedidas se enumeran en la
Secci\'on~\ref{subsec:leakage}, donde cada uno se empareja con la evidencia
de auditor\'ia que lo establece como riesgo concreto y no hipot\'etico. El
mecanismo es una lista negra expl\'icita m\'as una prueba automatizada que
hace fallar la construcci\'on si una columna prohibida alcanza la matriz de
caracter\'isticas, que es la forma operativa de la separaci\'on
aprender--predecir~\cite{kaufman_leakage_2011}.

\subsection{Plan de \emph{baselines}}
\label{subsec:baselines}
Cada pelda\~no de la escalera se justifica por lo que el anterior dej\'o sin
explicar.
\begin{enumerate}
  \item \textbf{B0 --- \emph{baseline} trivial.} Un clasificador de
        probabilidad a priori. Fija el suelo; lo que quede por debajo es
        ruido.
  \item \textbf{B1 --- heur\'istica de dominio.} Una regla \'unica sobre una
        diferencia de recursos, definida tras el an\'alisis exploratorio.
        Responde si hace falta aprendizaje autom\'atico en absoluto.
  \item \textbf{B2 --- regresi\'on log\'istica regularizada.}
        \textquestiondown Es la relaci\'on esencialmente lineal? Aporta coeficientes
        inspeccionables y probabilidades que pueden calibrarse. Es adem\'as
        el modelo que la evidencia revisada trata como referencia honesta:
        fue el algoritmo m\'as com\'un en los quince estudios cl\'inicos de
        cambio temporal~\cite{guo_temporalshift_2021}, qued\'o a 0.016 de
        acierto del mejor modelo en SC2EGSet~\cite{bialecki_sc2egset_2023},
        y estuvo entre los menos inflados por un protocolo defectuoso en el
        banco de pruebas de
        validaci\'on~\cite{karbalaie_participantaware_2026}.
  \item \textbf{B3 --- \'arbol de decisi\'on y \emph{random forest}.}
        \textquestiondown Existen interacciones y no linealidades que el modelo lineal
        pierde?
  \item \textbf{B4 --- \emph{gradient boosting}} (XGBoost, LightGBM o
        CatBoost), incluido solo si B3 muestra que el margen lo justifica.
        \textquestiondown Cu\'anto compra la complejidad, y a qu\'e costo de
        interpretabilidad? La evidencia revisada fija una expectativa baja
        sobre la respuesta: la mejor soluci\'on de \emph{Hearthstone}
        super\'o al \emph{baseline} de una red de dos capas por menos de un
        dos por ciento de AUC~\cite{janusz_hearthstone_2017}.
  \item \textbf{\emph{Baseline} de literatura.} La banda de $0.78$--$0.80$
        de AUC declarada en la Secci\'on~\ref{subsec:litbaseline}, con sus
        l\'imites declarados de comparabilidad.
\end{enumerate}

Las ablaciones por familia de caracter\'isticas, progreso, recursos,
tablero, energ\'ia y contexto de mazo, se ejecutan bajo el mismo
protocolo, de modo que la contribuci\'on marginal de cada bloque se mida en
lugar de afirmarse. Una ganancia grande del bloque de contexto de mazo se
interpreta como evidencia de memorizaci\'on del metajuego y no de
comprensi\'on, que es la lectura que Bertram \textit{et al.} justifican: una
representaci\'on que identifica el contenido en lugar de describirlo puede
ser excelente dentro de la distribuci\'on e in\'util fuera de
ella~\cite{bertram_cardrepresentations_2024}.

\subsection{Plan de explicabilidad}
\label{subsec:explainability}
La contribuci\'on no puede ser una puntuaci\'on de titular; la pregunta
sustantiva es cu\'ales caracter\'isticas del estado transportan informaci\'on
sobre la victoria. Se proponen cuatro instrumentos, cada uno ligado a una
pregunta. Los coeficientes de la regresi\'on log\'istica responden si la
asociaci\'on de una caracter\'istica con la victoria tiene la direcci\'on que
el conocimiento del dominio predice. La importancia por permutaci\'on
agrupada, permutando dentro de los episodios, nunca entre ellos,
responde cu\'anto aporta una caracter\'istica una vez respetada la
dependencia; Karbalaie \textit{et al.} se\~nalan expl\'icitamente que las
estimaciones de importancia dependen del dise\~no de validaci\'on y pueden
ser inestables bajo particiones a nivel de
fila~\cite{karbalaie_participantaware_2026}, raz\'on por la cual la
agrupaci\'on forma parte de la medici\'on y no es un detalle de ella. Las
ablaciones por familia responden cu\'anto vale cada bloque de
caracter\'isticas. Los valores SHAP se emplean \'unicamente cuando el modelo
elegido es un \emph{ensemble} de \'arboles y la pregunta es la atribuci\'on
por instancia, por qu\'e \emph{esta} posici\'on recibi\'o esta
puntuaci\'on, puesto que esa es la pregunta que los tres instrumentos
anteriores no pueden responder; no se incluyen meramente porque est\'en
disponibles.

\subsection{Contrato de reproducibilidad}
\label{subsec:reproducibility}
Las semillas se fijan y se versionan; la auditor\'ia ya se ejecuta bajo una
\'unica semilla declarada. Las configuraciones viven en ficheros y no en
banderas de l\'inea de comandos. La representaci\'on tabular se materializa
una sola vez en formato columnar, puesto que el cuello de botella
computacional es el an\'alisis sint\'actico y no el entrenamiento. Los datos
crudos nunca se versionan; el repositorio lleva los \emph{scripts} y las
instrucciones para obtenerlos. Esto refleja la pr\'actica del descriptor
SC2EGSet, que publica sus herramientas de extracci\'on, sus registros de
procesamiento y sus estad\'isticas resumidas por paquete de repeticiones
junto a los datos~\cite{bialecki_sc2egset_2023}, y atiende una debilidad
documentada de la literatura sobre cambio temporal, donde el c\'odigo estaba
disponible en solo dos de quince estudios
revisados~\cite{guo_temporalshift_2021}.

\subsection{Riesgos y mitigaci\'on}
\label{subsec:risks}
\begin{table}[!t]
\centering
\caption{Riesgos identificados y su mitigaci\'on.}
\label{tab:risks}
\footnotesize
\begin{tabularx}{\columnwidth}{@{}P{1.9cm} c Y@{}}
\toprule
\textbf{Riesgo} & \textbf{Impacto} & \textbf{Mitigaci\'on / disparador}\\
\midrule
Instant\'aneas de un mismo episodio repartidas entre particiones & A &
Agrupar por \code{episode\_id} en todas partes; prueba automatizada que
verifica intersecci\'on vac\'ia entre particiones\\
Un campo prohibido alcanza el vector de caracter\'isticas & A &
Lista negra expl\'icita; la construcci\'on falla si una columna prohibida
est\'a presente\\
Fuga introducida por decisiones de dise\~no & M &
Reglas de exclusi\'on y ventanas de turno declaradas antes de tocar el
conjunto de prueba~\cite{kaufman_leakage_2011}\\
Preprocesamiento ajustado sobre el conjunto completo & M &
Todas las transformaciones se ajustan dentro del pliegue de
entrenamiento~\cite{sasse_leakage_2025}\\
Cambio de distribuci\'on entre d\'ias tempranos y tard\'ios & M &
Reportar el protocolo temporal por separado y declarar cualquier ca\'ida
como resultado, no como fallo~\cite{guo_temporalshift_2021}\\
Memorizaci\'on de mazos en lugar de comprensi\'on del estado & M &
Protocolo de reserva de mazos y ablaci\'on del bloque de contexto de
mazo~\cite{bertram_cardrepresentations_2024}\\
Caracter\'isticas trivialmente predictivas cerca del final & M &
An\'alisis por ventana de turno; las conclusiones tempranas se sostienen
solo sobre la ventana temprana\\
Versi\'on del motor confundida con el tiempo de calendario & M &
Registrar la versi\'on como covariable de auditor\'ia y declarar la
confusi\'on\\
Lectura optimista de una puntuaci\'on no calibrada & M &
Reportar una regla de puntuaci\'on propia y curvas de fiabilidad junto al
AUC~\cite{guo_temporalshift_2021}\\
\bottomrule
\end{tabularx}
\end{table}

\subsection{Integraci\'on prevista con el agente}
\label{subsec:agent}
En la Etapa 3 el estimador entrenado es consumido por un agente inteligente
como evaluador de posiciones. La interfaz es deliberadamente estrecha: el
agente entrega un estado, el mismo extractor de caracter\'isticas usado en
entrenamiento produce $x$, y el modelo devuelve $f(x) \in [0,1]$. El agente
no entrena el modelo, no lo actualiza mediante interacci\'on y no lo
sustituye, de modo que el proyecto no se convierte en aprendizaje por
refuerzo en ninguna etapa. El consumo se verificar\'a mostrando que el
ordenamiento de posiciones candidatas por parte del agente cambia cuando el
estimador se reemplaza por el \emph{baseline} trivial B0, es decir, que el
agente est\'a efectivamente usando la predicci\'on. La calibraci\'on importa
precisamente por esta raz\'on: un evaluador que ordena correctamente las
posiciones pero est\'a sistem\'aticamente sobreconfiado enga\~nar\'a a
cualquier regla de decisi\'on que compare sus salidas contra un umbral fijo.

% ===========================================================================
% 6.2.6  ANALISIS TECNICO DEL TRACK A
% ===========================================================================
\section{An\'alisis T\'ecnico del Track Seleccionado}
\label{sec:trackanalysis}

El track aprobado es el Track A, aprendizaje autom\'atico cl\'asico sobre
datos tabulares. El an\'alisis que sigue se apoya en una auditor\'ia del
corpus realizada antes de cualquier modelado; produjo un registro por
episodio para cada episodio auditado, un censo de campos a nivel de
instant\'anea y un conjunto de compuertas de decisi\'on, todo reproducible
bajo una \'unica semilla fija.

\emph{Regla epist\'emica de esta secci\'on.} Toda cifra que aparece a
continuaci\'on fue medida por nuestra propia auditor\'ia del corpus y es
reproducible a partir de los \emph{scripts} de auditor\'ia. Las magnitudes
que no han sido medidas se nombran como pendientes, nunca se estiman.

\subsection{An\'alisis exploratorio}
\label{subsec:eda}

\subsubsection{Alcance de lo que ha sido medido}
El corpus publica 67 entregas diarias que abarcan del 2026-06-16 al
2026-08-21, declarando 329\,336 episodios y 1.29~TB sin comprimir. La
auditor\'ia cubre una muestra temporal de 10 de esos d\'ias, el
$14.9\%$ de los d\'ias y el $15.8\%$ de los episodios declarados, 
elegidos para abarcar las fases inicial, media y tard\'ia de la competencia,
con tres d\'ias consecutivos al comienzo de modo que pudiera medirse la
deriva d\'ia a d\'ia. Dentro de esa muestra se analizaron 52\,085 episodios
a partir de 182.6~GiB de JSON con cero errores de an\'alisis sint\'actico,
produciendo 7\,038\,378 instant\'aneas utilizables. Toda afirmaci\'on de
esta secci\'on est\'a acotada a esa muestra; nada de lo aqu\'i escrito
caracteriza a los 57 d\'ias no auditados.

\subsubsection{Tama\~no, estructura y unidad de observaci\'on}
La duraci\'on de las partidas est\'a fuertemente sesgada a la derecha. El
n\'umero de pasos por episodio tiene mediana 140 y media 137.1, con
cuartiles 101 y 171, un percentil 95 de 221 y un m\'aximo de 6\,900. El
recuento de turnos es mucho menor que el de pasos, mediana 13, media
14.3, percentil 95 de 27, m\'aximo 1\,035, porque un solo turno se
descompone en muchas acciones at\'omicas. Esa brecha gobierna c\'omo debe
definirse una ventana de turno: fijar el ``juego temprano'' en los turnos
uno a ocho ser\'ia una elecci\'on arbitraria contra una distribuci\'on cuya
mediana es 13, de modo que la ventana se define a partir de la
distribuci\'on emp\'irica y se declara antes de tocar el conjunto de prueba.

La unidad de observaci\'on es un estado fresco del asiento que act\'ua, y
esta es una decisi\'on medida antes que estil\'istica. En un sondeo de
139\,672 estados frescos, la observaci\'on perteneciente al asiento que no
act\'ua result\'o id\'entica a su propio estado anterior en 137\,989 casos,
es decir el $98.8\%$. Tratar cada asiento de cada paso como una fila
duplicar\'ia por tanto casi todas las observaciones sin a\~nadir
informaci\'on, y reponderar\'ia silenciosamente los episodios por su
duraci\'on. Restringirse al asiento activo da una raz\'on medida entre
estados activos y pasos de 0.9854, que es de donde procede la cifra de
7\,038\,378, y no de $2 \times n_{\text{pasos}}$.

El n\'umero de unidades \emph{independientes} es 52\,085. Siete millones de
filas y cincuenta y dos mil unidades independientes es el hecho m\'as
consecuente de esta secci\'on, y la Secci\'on~\ref{subsec:leakage} est\'a
construida sobre \'el.

\subsubsection{Distribuci\'on del objetivo a nivel de episodio}
El desenlace se deriva del campo de recompensas a nivel de episodio, que
resuelve el $99.86\%$ de los episodios. Las combinaciones observadas son
$(1,-1)$ en 27\,135 episodios, $(-1,1)$ en 24\,844, $(0,0)$ en 33, y una
recompensa que no es un n\'umero en 73. A nivel de episodio esto da un
$52.10\%$ de victorias para el asiento P0, un $47.70\%$ para P1, un
$0.063\%$ de empates y un $0.140\%$ de abandonos, sin episodios sin
resolver. Restringido a los 51\,979 episodios con ganador decidido, la tasa
de victoria de P0 es 0.5220 y la del jugador que mueve primero es 0.5287.

\subsubsection{Estructura temporal y deriva}
Los diez d\'ias auditados no son intercambiables. Medida con la divergencia
de Jensen--Shannon en base 2 entre el primer y el \'ultimo d\'ia auditado,
la distribuci\'on de mazos es esencialmente disjunta, con 1.0000, la
poblaci\'on de agentes casi tanto, con 0.9723, la distribuci\'on de cartas
se ha desplazado hasta 0.4889, mientras que la duraci\'on de las partidas
apenas se ha movido, con 0.1221. La diversidad no decae de forma suave: los
mazos distintos por d\'ia van de 1\,268 el 2026-06-23 a 48 el 2026-07-17, y
los equipos distintos de 2\,309 a 79 en esos mismos dos d\'ias. La mediana
de la puntuaci\'on media de los episodios sube de 628 el primer d\'ia a
1\,141 el 2026-07-17. Seis versiones del motor coexisten en la muestra, de
la \code{1.30.1} a la \code{1.32.4}, y est\'an ordenadas en el tiempo, lo
que significa que tiempo de calendario y versi\'on del motor est\'an
confundidos en la muestra auditada. La auditor\'ia no puede separar la
deriva del metajuego del cambio del motor, y este trabajo reportar\'a esa
confusi\'on en lugar de resolverla.

Un hallazgo de deriva constituye una amenaza directa a una caracter\'istica
concreta. La probabilidad de que el jugador que mueve primero sea el asiento
P0 era 0.5089 el 2026-06-29 y es 0.9991 desde el 2026-07-17 en adelante. En
los d\'ias tard\'ios, identidad de asiento y orden de turno son casi la
misma variable.

\subsection{Diccionario de datos}
\label{subsec:datadict}

Un censo de campos sobre 139\,672 estados frescos encontr\'o 34 campos de
estado presentes en el $100\%$ de las instant\'aneas inspeccionadas: 11
campos globales, 13 campos por jugador y 10 campos por Pok\'emon sobre
1\,118\,776 instancias de Pok\'emon. La Tabla~\ref{tab:datadict} agrupa las
caracter\'isticas candidatas por familia y tipo. El diccionario es
preliminar y est\'a sujeto a revisi\'on tras el an\'alisis exploratorio de
la tabla materializada.

\begin{table}[!t]
\centering
\caption{Diccionario de datos por familia de caracter\'isticas. La
``presencia'' es la proporci\'on medida de instant\'aneas inspeccionadas en
las que el campo de origen existe.}
\label{tab:datadict}
\footnotesize
\begin{tabularx}{\columnwidth}{@{}P{1.5cm} P{1.35cm} Y@{}}
\toprule
\textbf{Familia} & \textbf{Tipo} & \textbf{Campos, derivaci\'on y presencia}\\
\midrule
Progreso & entero, binario &
\code{turn} y \code{turnActionCount} directos; el indicador de jugador
inicial derivado de \code{firstPlayer} y \code{yourIndex};
\code{energyAttached} y \code{retreated} binarios. Presencia $100\%$.\\
Recursos & entero &
\code{deckCount} y \code{handCount} directos; el tama\~no de la mano
derivado de la lista de la mano propia; el tama\~no del descarte derivado.
Presencia $100\%$.\\
Premios & entero &
Ranuras de premio restantes, derivadas contando las ranuras no vac\'ias.
Presencia $100\%$; el contenido est\'a siempre oculto.\\
Tablero & entero, continuo, binario &
Tama\~no de la banca frente a \code{benchMax}; \code{hp}, \code{maxHp} y su
raz\'on; totales de puntos de vida visibles; recuento de Pok\'emon
da\~nados; \code{energies}, \code{energyCards}, \code{tools};
\code{preEvolution} y \code{appearThisTurn} como binarios. Presencia
$100\%$.\\
Estado alterado & binario &
\code{poisoned}, \code{burned}, \code{asleep}, \code{paralyzed} y
\code{confused} por jugador; \code{stadium}, \code{stadiumPlayed} y
\code{supporterPlayed} globales. Presencia $100\%$.\\
Diferenciales & entero, continuo &
Propio menos rival para recuentos de mano, premios, banca, puntos de vida,
energ\'ia y mazo. Derivados.\\
Contexto de mazo & categ\'orico &
Huella del mazo sobre un cat\'alogo de 2\,376 mazos distintos. Alta
cardinalidad; se admite \'unicamente en la ablaci\'on de contexto de mazo.\\
\midrule
Agrupaci\'on \newline (nunca caracter\'isticas) & --- &
\code{episode\_id}, \code{date} y el \'indice de turno. Se arrastran para
agrupaci\'on, estratificaci\'on y reporte por ventana.\\
\bottomrule
\end{tabularx}
\end{table}

Dos campos se excluyen sobre bases medidas antes de cualquier modelado. El
campo de resultado del estado actual contiene el valor $-1$ en el
$100\%$ de los 139\,672 estados frescos inspeccionados: es constante y no
transporta informaci\'on en el momento de la predicci\'on. El campo
\code{looking} es nulo en 137\,029 de esos estados, es decir el
$98.1\%$, y su nulidad es una propiedad del juego y no un defecto,
v\'ease la Secci\'on~\ref{subsec:imputation}.

\subsection{Calidad de los datos}
\label{subsec:quality}

\subsubsection{Qu\'e ha sido verificado}
Los 52\,085 episodios auditados se analizaron sin error. Una \'unica firma
de esquema de nivel superior cubre el $100\%$ de ellos y
\code{schema\_version} vale 1 en todos, de modo que el corpus es
estructuralmente homog\'eneo pese a las seis versiones del motor. Los
\'indices de turno fueron mon\'otonos en todos los episodios inspeccionados,
con cero violaciones. Todos los episodios auditados aparecen en el manifiesto
de su d\'ia; a la inversa, cinco episodios declarados en el manifiesto del
2026-08-04 est\'an ausentes de los ficheros, de 4\,816 declarados.

\subsubsection{Duplicados}
Tres comprobaciones independientes de duplicidad devolvieron cero: ning\'un
\code{episode\_id} repetido, ning\'un SHA-256 repetido sobre los bytes del
fichero, y ninguna huella can\'onica repetida, esta \'ultima calculada sobre
la semilla de configuraci\'on, ambos mazos de sesenta cartas en orden, la
secuencia completa de acciones de ambos asientos y las recompensas,
excluyendo deliberadamente los metadatos vol\'atiles. Ning\'un
\code{episode\_id} aparece en m\'as de un d\'ia. El juego contra uno mismo,
donde el mismo equipo ocupa ambos asientos, ocurre en cero de los 52\,085
episodios.

La duplicaci\'on exacta no es, por tanto, una preocupaci\'on a nivel de
episodio. La duplicaci\'on funcional dentro de un episodio es una cuesti\'on
distinta y sigue abierta: instant\'aneas consecutivas de la misma partida
difieren en una acci\'on at\'omica y son casi id\'enticas por
construcci\'on. En la muestra sondeada, los 139\,672 estados frescos fueron
distintos dentro de sus turnos, de modo que no hay dos filas id\'enticas a
nivel de bytes, pero las cuasi duplicadas son la norma antes que la
excepci\'on. Esto no es un defecto que deba eliminarse, es la estructura
de dependencia de los datos, y se aborda mediante agrupaci\'on y no mediante
deduplicaci\'on.

\subsubsection{Comprobaciones a\'un pendientes sobre la representaci\'on tabular}
La auditor\'ia verific\'o el corpus crudo. Lo siguiente debe reverificarse
sobre la tabla materializada, y se consigna aqu\'i como pendiente antes que
como supuesto: la ausencia de valores por columna; los rangos imposibles,
como puntos de vida negativos, puntos de vida por encima del m\'aximo,
recuentos negativos o tama\~no de banca por encima de \code{benchMax}; la
consistencia l\'ogica entre campos relacionados; las filas duplicadas dentro
y entre episodios; y la cardinalidad efectiva de las columnas categ\'oricas.

\subsection{Valores ausentes e imputaci\'on}
\label{subsec:imputation}

La pol\'itica distingue dos mecanismos, porque tratarlos por igual ser\'ia
un error.

\textbf{Ausencia estructural.} Un campo es nulo porque el juego lo define
as\'i: no hay Pok\'emon activo, no hay estadio en juego, la carta est\'a
oculta. El caso m\'as claro es \code{looking}, nulo en el $98.1\%$ de los
estados inspeccionados, y el segundo es la mano del rival, nula en el
$100\%$ de ellos. No son valores estad\'isticamente ausentes y no se
imputan. Se codifican sem\'anticamente, con un indicador binario de presencia
m\'as un centinela expl\'icito. Imputar una media aqu\'i fabricar\'ia un
estado de juego que no puede ocurrir, y adem\'as destruir\'ia la se\~nal
misma que el campo transporta: el hecho de que no se est\'e mirando nada es
informaci\'on.

\textbf{Ausencia inesperada.} Un campo requerido est\'a ausente porque un
episodio est\'a malformado o el analizador sint\'actico es incompleto. El
procedimiento es medir primero la frecuencia; si es residual, excluir la
observaci\'on o el episodio dejando constancia del criterio y del recuento;
si es sistem\'atica, tratarla como un defecto del analizador o un cambio de
esquema y corregir la tuber\'ia en lugar de parchear los datos. Ninguna
variable num\'erica se imputa autom\'aticamente antes de haber estudiado su
mecanismo de ausencia.

Sea cual sea el mecanismo elegido, el estimador se ajusta dentro del pliegue
de entrenamiento. Sasse \textit{et al.} nombran expl\'icitamente la
imputaci\'on entre los pasos de preprocesamiento que producen fuga cuando se
estiman sobre el conjunto completo antes de
particionar~\cite{sasse_leakage_2025}. La estrategia definitiva se fija tras
el an\'alisis exploratorio de la tabla materializada y se documenta con la
frecuencia que la justifica.

\subsection{Detecci\'on de valores at\'ipicos}
\label{subsec:outliers}

En este corpus un valor extremo suele ser un at\'ipico v\'alido antes que un
error. Los pasos por episodio alcanzan 6\,900 frente a una mediana de 140, y
el recuento de turnos alcanza 1\,035 frente a una mediana de 13 y un
percentil 95 de 27. Una partida de 1\,035 turnos merece investigaci\'on, no
borrado: puede ser un tiempo agotado, un bucle de acciones, una interacci\'on
inusual de cartas, un fallo del simulador o una partida genuinamente larga, y
esos cinco casos exigen respuestas distintas.

El procedimiento propuesto es, por tanto, diagn\'ostico antes que
correctivo. Los candidatos se marcan por distancia distribucional, una
regla intercuartil o de puntuaci\'on $z$ robusta sobre las caracter\'isticas
de duraci\'on y recuento, y por rangos imposibles en el dominio, que es
una prueba distinta con un significado distinto. Los episodios marcados se
inspeccionan caso por caso contra los estados registrados para ellos. Solo se
eliminan los registros demostrablemente corruptos, y cualquier regla de
exclusi\'on se declara antes de examinar la partici\'on de prueba; una regla
elegida despu\'es de ver su efecto sobre la puntuaci\'on de prueba es una
fuga por decisi\'on de dise\~no en el sentido de Kaufman \textit{et
al.}~\cite{kaufman_leakage_2011}. Se\~nalamos un contraste con la
pr\'actica seguida en un corpus comparable: los autores de SC2EGSet
reemplazaron por la mediana los valores que exced\'ian tres desviaciones
est\'andar en sus propios experimentos~\cite{bialecki_sc2egset_2023}.
Preferimos no adoptar esa regla por defecto aqu\'i, porque en un juego por
turnos la cola de la distribuci\'on de duraci\'on puede ser exactamente
donde vive el comportamiento interesante.

\subsection{Desbalance de clases}
\label{subsec:imbalance}

A nivel de episodio el objetivo est\'a cerca del equilibrio: $52.10\%$
frente a $47.70\%$, con empates en $0.063\%$ y abandonos en
$0.140\%$.

Esa no es, sin embargo, la distribuci\'on que ver\'a el modelo. La etiqueta
se define desde la perspectiva del jugador que act\'ua, y cada partida aporta
un n\'umero distinto de instant\'aneas desde cada bando, de modo que la
prevalencia de clases de la tabla final es funci\'on de la tabularizaci\'on y
\textbf{a\'un no ha sido medida}. Debe recalcularse sobre la tabla
materializada antes de tomar decisi\'on alguna, y esta secci\'on se
completar\'a con esa medici\'on.

La regla de decisi\'on queda no obstante fijada de antemano, de modo que no
pueda elegirse para acomodar un resultado. Si el desbalance resultante es
leve no se corregir\'a: corregir un desbalance que no necesita correcci\'on
distorsiona las probabilidades predichas, y esas probabilidades son la salida
de este proyecto y no un subproducto. Si es apreciable, se preferir\'a la
ponderaci\'on de clases al remuestreo, por la misma raz\'on. No se usar\'a
sobremuestreo sint\'etico sin justificaci\'on expl\'icita, y bajo ninguna
circunstancia fuera del pliegue de entrenamiento, Sasse \textit{et al.}
enumeran el sobremuestreo aplicado antes de particionar entre sus casos
documentados de fuga de prueba a entrenamiento~\cite{sasse_leakage_2025}.

Los 33 empates y los 73 abandonos se tratan por separado. Los empates se
excluyen del objetivo binario y se documentan; su volumen no puede sostener
una tercera clase. Los abandonos tienen ganador identificable, pero el final
no procede de la din\'amica del juego, los estados registrados son tiempos
agotados, errores y estados inv\'alidos, de modo que se excluyen del
entrenamiento y se reservan para un an\'alisis de sensibilidad.

\subsection{Dependencia entre muestras y prevenci\'on de la fuga de datos}
\label{subsec:leakage}

Esta es la subsecci\'on de la que depende la validez de todo lo dem\'as, y
el enunciado del curso penaliza la fuga confirmada m\'as severamente que
cualquier otro fallo aislado.

\subsubsection{La estructura de dependencia, medida}
Cada episodio aporta en promedio unas 135 instant\'aneas, e instant\'aneas
consecutivas difieren en una \'unica acci\'on at\'omica. Las filas son por
tanto fuertemente dependientes dentro de un episodio e independientes solo
entre episodios. Medimos la consecuencia directamente: bajo una partici\'on
aleatoria $70/30$ a nivel de instant\'anea, la fracci\'on esperada de
episodios presentes en \emph{ambas} particiones es 0.999945. Dicho de otro
modo, el protocolo ingenuo contamina esencialmente todo el corpus.

La literatura dice qu\'e le har\'ia eso a las cifras. Karbalaie
\textit{et al.} midieron un acierto de 0.91 para \emph{extra trees} bajo
validaci\'on cruzada estratificada de diez pliegues frente a 0.66 dejando un
participante fuera sobre los mismos datos, e identificaron la causa en que el
tama\~no muestral efectivo es el n\'umero de participantes y no el de
ensayos~\cite{karbalaie_participantaware_2026}. Sasse \textit{et al.}
clasifican precisamente este caso, no tener en cuenta muestras
dependientes, como v\'ia primaria de fuga de prueba a entrenamiento y lo
demuestran sobre pares de hermanos en una cohorte de
neuroimagen~\cite{sasse_leakage_2025}. Nuestro corpus es una instancia
m\'as extrema de la misma estructura: no dos muestras relacionadas, sino
135.

\subsubsection{El momento de la predicci\'on}
Un campo puede existir en el JSON y aun as\'i ser inv\'alido como
caracter\'istica. La prueba que gobierna es el requisito de la no m\'aquina
del tiempo de Kaufman \textit{et al.}: un modelo leg\'itimo solo puede usar
informaci\'on de un instante no posterior al del
objetivo~\cite{kaufman_leakage_2011}, afinada por la pregunta que Sasse
\textit{et al.} recomiendan hacer de cada columna, si ese campo
estar\'ia disponible en el despliegue~\cite{sasse_leakage_2025}. Aplicar esa
prueba campo por campo produce cuatro clases.

\emph{Seguras.} Los campos de estado presentes en el $100\%$ de las
instant\'aneas, le\'idos de la observaci\'on propia del asiento activo en el
turno $\tau$, junto con los recuentos, razones y diferenciales derivados de
ellos. Son lo que el propio agente vio cuando decidi\'o.

\emph{Requieren investigaci\'on.} El indicador de jugador inicial, porque
asiento y orden de turno est\'an casi perfectamente confundidos desde el
2026-07-17 en adelante y el modelo podr\'ia aprender el asiento en lugar del
tempo; y el registro acumulado de acciones, que es leg\'itimo \'unicamente si
se trunca estrictamente antes de $\tau$.

\emph{Fuga directa, excluidas por construcci\'on.} El campo de recompensas a
nivel de episodio y el resultado final; cualquier recompensa por paso, que es
no nula \'unicamente en el \'ultimo paso; la carga \'util de visualizaci\'on,
presente en el $100\%$ de los episodios, que expone el mazo restante de
cada jugador en orden, informaci\'on perfecta que nunca aparece en la
observaci\'on y que el agente no ten\'ia; las marcas de tiempo y los
identificadores derivados del fichero; y los nombres de equipo o de agente,
que act\'uan como sustituto de la habilidad. Lo \'ultimo no es hipot\'etico:
la mediana de la puntuaci\'on media de los episodios seleccionados cada
d\'ia sube de 628 a m\'as de 1\,100 a lo largo de la muestra, de modo que un
identificador de agente codifica parcialmente el desenlace.

\emph{Excluidas por constantes.} El campo de resultado del estado actual,
fijo en $-1$ en todos los estados frescos.

\subsubsection{Informaci\'on oculta, verificada antes que supuesta}
El entorno ya retiene lo que debe retener, y lo comprobamos en lugar de
confiar. A lo largo de 139\,672 estados inspeccionados, la mano del rival fue
nula en el $100\%$ de los casos con cero fugas detectadas, mientras que el
recuento de cartas del rival estuvo siempre disponible como entero. Las
1\,281\,678 ranuras de premio inspeccionadas, propias y del rival por
igual, fueron todas nulas. La fuga de informaci\'on oculta no es, por
tanto, un riesgo que debamos mitigar; es una propiedad que no debemos
destruir, raz\'on por la cual las caracter\'isticas se extraen exclusivamente
de la observaci\'on propia del asiento y nunca de la carga \'util de
visualizaci\'on.

\subsubsection{Contramedidas}
La Tabla~\ref{tab:leakage} empareja cada vector con la evidencia que lo
convierte en un riesgo concreto y con el control aplicado. Cada control se
impone mediante una prueba automatizada y no por convenci\'on, que es la
forma operativa de la separaci\'on aprender--predecir de Kaufman \textit{et
al.}~\cite{kaufman_leakage_2011}.

\begin{table}[!t]
\centering
\caption{Vectores de fuga de datos, la evidencia de auditor\'ia que
establece cada uno como riesgo real, y el control aplicado.}
\label{tab:leakage}
\footnotesize
\begin{tabularx}{\columnwidth}{@{}P{1.55cm} P{2.35cm} Y@{}}
\toprule
\textbf{Vector} & \textbf{Evidencia medida} & \textbf{Control}\\
\midrule
Fuga de episodio &
Una partici\'on aleatoria $70/30$ por filas sit\'ua 0.999945 de los
episodios en ambas particiones &
Partici\'on y pliegues agrupados por \code{episode\_id}; prueba que verifica
intersecci\'on vac\'ia\\
Informaci\'on futura &
Carga \'util de visualizaci\'on presente en el $100\%$ de los episodios,
que expone el mazo restante ordenado &
Caracter\'isticas extra\'idas solo de la observaci\'on del asiento en
$\tau$; carga \'util en lista negra\\
Sustituto de la habilidad &
La mediana de la puntuaci\'on media sube de 628 a m\'as de 1\,100 a lo largo
de la muestra &
Metadatos de agente y equipo usados solo para auditor\'ia y agrupaci\'on,
nunca como caracter\'isticas\\
Memorizaci\'on de mazos &
2\,376 mazos, pero los 10 primeros cubren el $39.94\%$ de las
instancias; un d\'ia tiene solo 48 &
Protocolo de reserva de mazos; contexto de mazo confinado a una \'unica
ablaci\'on declarada\\
Informaci\'on oculta &
$0$ fugas en 139\,672 estados; 1\,281\,678 ranuras de premio todas nulas &
Analizador que lee solo la observaci\'on propia, con una aserci\'on
automatizada de privacidad\\
Metadatos &
52\,085 \code{episode\_id} distintos sobre 52\,085 filas &
Columnas de agrupaci\'on excluidas de la matriz de caracter\'isticas por
construcci\'on\\
Preprocesamiento &
--- (riesgo de dise\~no, no propiedad del corpus) &
Todas las transformaciones se ajustan dentro del pliegue de
entrenamiento~\cite{sasse_leakage_2025}\\
Temporal &
Divergencia de Jensen--Shannon de mazos de 1.0000 entre el primer y el
\'ultimo d\'ia auditado &
El protocolo temporal corta por fecha, nunca al
azar~\cite{kaufman_leakage_2011}\\
\bottomrule
\end{tabularx}
\end{table}

\subsection{Representatividad}
\label{subsec:representativeness}

El corpus no es una muestra aleatoria del juego de \emph{Pok\'emon TCG}, y
el editor lo dice: las repeticiones se seleccionan a diario, ordenadas por
puntuaci\'on media de los agentes, y se truncan a 20~GiB por
d\'ia~\cite{ptcg_dataset_2026}. De ah\'i se siguen tres consecuencias, y las
tres est\'an medidas. La poblaci\'on es de agentes artificiales, no de
humanos. Es la franja alta de la escalera de clasificaci\'on, con la mediana
de la puntuaci\'on media ascendiendo a lo largo de la muestra. Y est\'a
concentrada: los diez mazos m\'as frecuentes suman el $39.94\%$ de las
instancias de mazo y los cien m\'as frecuentes el $74.76\%$, con d\'ias
individuales que caen a 48 mazos distintos.

Esto refleja la limitaci\'on que los autores de SC2EGSet declaran para su
propio corpus, que solo contiene repeticiones de nivel de torneo y por tanto
no admite an\'alisis entre niveles de habilidad~\cite{bialecki_sc2egset_2023},
y es una versi\'on m\'as aguda del problema de poblaci\'on al que Janusz
\textit{et al.} se enfrentan desde la direcci\'on opuesta, pues sus estados
fueron generados por agentes que act\'uan al
azar~\cite{janusz_hearthstone_2017}. Le\'idos en conjunto, ambos acotan
nuestro caso: un corpus procede de la poblaci\'on m\'as d\'ebil posible y el
otro de la m\'as fuerte, y ninguno es la poblaci\'on general. Cualquier
afirmaci\'on de este proyecto es, por tanto, una afirmaci\'on sobre agentes
artificiales competitivos de una franja de puntuaci\'on espec\'ifica durante
una muestra espec\'ifica de diez d\'ias, y la Secci\'on~\ref{sec:ethics}
trata la tentaci\'on de enunciarla de forma m\'as amplia como el principal
riesgo \'etico.

% ===========================================================================
% 6.2.7  CONSIDERACIONES ETICAS PRELIMINARES
% ===========================================================================
\section{Consideraciones \'Eticas Preliminares}
\label{sec:ethics}

\subsection{Sensibilidad de los datos}
El corpus no contiene datos cl\'inicos, biom\'etricos ni conductuales sobre
individuos identificados: sus registros son partidas simuladas entre agentes
de software, publicadas bajo una licencia
CC0-1.0~\cite{ptcg_dataset_2026}. Cabe una salvedad. Cada episodio lleva los
nombres de equipo y de agente de los competidores, y en la muestra auditada
esos nombres son seud\'onimos elegidos por sus due\~nos que en algunos casos
se asemejan a nombres personales. Identifican a participantes de una
competencia p\'ublica y no a sujetos de estudio, pero siguen siendo
atribuibles a personas. Los tratamos por tanto como metadatos y no como
material: se usan para auditor\'ia, agrupaci\'on y la comprobaci\'on de juego
contra uno mismo, quedan excluidos de la matriz de caracter\'isticas por la
lista negra descrita en la Secci\'on~\ref{subsec:leakage}, y no se
publicar\'a ning\'un ranking ni comparaci\'on agregada de competidores
nombrados. Esto coincide con el control de fuga, lo cual es conveniente pero
no es la raz\'on de la decisi\'on.

\subsection{Sesgo y representatividad}
El principal riesgo de este proyecto no es el da\~no a un sujeto de datos;
es la sobreafirmaci\'on. Los episodios se seleccionan a diario por
puntuaci\'on media de los agentes y se truncan a un volumen diario
fijo~\cite{ptcg_dataset_2026}, de modo que el corpus es la franja alta de una
escalera de agentes artificiales y no una muestra del juego de
\emph{Pok\'emon TCG}. La auditor\'ia cuantifica la concentraci\'on que se
sigue de ello: los diez mazos m\'as frecuentes cubren el $39.94\%$ de las
instancias de mazo, un d\'ia auditado contiene solo 48 mazos distintos, y el
conjunto de mazos del \'ultimo d\'ia auditado es esencialmente disjunto del
del primero. Un modelo ajustado aqu\'i aprende un metajuego estrecho y en
movimiento.

De ah\'i se siguen tres lecturas err\'oneas, y cada una merece ser nombrada
porque cada una es f\'acil de cometer. Extrapolar a jugadores humanos
carecer\'ia de respaldo: no aparece ninguna partida humana en el corpus, y
los dos corpus comparables de deportes electr\'onicos que revisamos declaran
exactamente esta limitaci\'on para s\'i
mismos~\cite{bialecki_sc2egset_2023,janusz_hearthstone_2017}. Extrapolar al
juego completo tampoco tendr\'ia respaldo, puesto que un modelo que nunca vio
un arquetipo de mazo carece de base para juzgar posiciones construidas en
torno a \'el, Bertram \textit{et al.} muestran c\'omo se ve ese fallo
cuando el contenido se codifica por identidad en lugar de por
descripci\'on~\cite{bertram_cardrepresentations_2024}. Y tratar la salida
como una propiedad estable de una posici\'on tampoco tendr\'ia respaldo a lo
largo del tiempo, porque la poblaci\'on cambia de forma medible en diez
d\'ias.

\subsection{Interpretaci\'on de la probabilidad}
Una estimaci\'on de probabilidad de victoria invita a un mal uso
espec\'ifico: leer un n\'umero entre cero y uno como una posibilidad bien
fundada cuando puede ser solo una puntuaci\'on bien ordenada. La distinci\'on
no es pedante. Guo \textit{et al.} encontraron deterioro de la calibraci\'on
en todos los estudios que la midieron bajo cambio temporal, frente a
deterioro de la discriminaci\'on en solo tres de
doce~\cite{guo_temporalshift_2021}, lo que significa que un modelo puede
seguir ordenando posiciones correctamente mucho despu\'es de que sus
probabilidades declaradas hayan dejado de ser ciertas. Ninguno de los dos
estudios comparables sobre estados de juego reporta calibraci\'on en
absoluto~\cite{janusz_hearthstone_2017,bialecki_sc2egset_2023}. Reportar una
regla de puntuaci\'on propia y curvas de fiabilidad junto al AUC es por
tanto un compromiso \'etico tanto como metodol\'ogico: es lo que permite a
un lector saber cu\'al de las dos cosas es nuestro n\'umero.

\subsection{Mitigaciones previstas antes del entrenamiento}
\begin{itemize}
  \item Toda cifra reportada lleva su alcance: los diez d\'ias auditados, la
        franja de puntuaci\'on y las versiones del motor presentes.
  \item La discriminaci\'on y la calibraci\'on se reportan por separado, y
        los resultados bajo los protocolos temporal y de reserva de mazos se
        reportan aparte del protocolo primario en lugar de promediarse con
        \'el.
  \item Los identificadores de agente y equipo quedan excluidos de la matriz
        de caracter\'isticas, y no se publican resultados por competidor.
  \item Los intervalos de confianza se calculan sobre episodios, de modo que
        la incertidumbre no se subestime contando filas correlacionadas como
        evidencia independiente.
  \item El resumen, las conclusiones y cualquier s\'intesis p\'ublica
        declaran la poblaci\'on sobre la que se ajust\'o el modelo, de modo
        que la limitaci\'on viaje con el resultado en lugar de quedar
        confinada a esta secci\'on.
\end{itemize}

\subsection{Diferido a la Etapa 3}
La \'etica del despliegue, los escenarios de mal uso y los riesgos
operativos de un agente que consuma este estimador, en particular el
efecto de un evaluador sobreconfiado sobre las decisiones del agente, se
abordan en la etapa final, cuando exista un agente.

% ===========================================================================
% DECLARACION DE USO DE HERRAMIENTAS DE IA
% ===========================================================================
\section{Declaraci\'on de Uso de Herramientas de Inteligencia Artificial}
\label{sec:aideclaration}

\begin{table}[!t]
\centering
\caption{Uso de herramientas de inteligencia artificial en esta versi\'on
individual de la Etapa 1.}
\label{tab:aiusage}
\footnotesize
\begin{tabularx}{\columnwidth}{@{}P{1.5cm} P{1.5cm} Y@{}}
\toprule
\textbf{Integrante} & \textbf{Herramienta} & \textbf{D\'onde y para qu\'e}\\
\midrule
S. Hern\'andez &
Claude (Claude Code, Opus~5) &
Organizaci\'on y cribado de la literatura; recuperaci\'on de textos
completos desde fuentes de los editores y de acceso abierto; extracci\'on de
figuras, tablas y limitaciones declaradas de los siete trabajos revisados;
resoluci\'on de metadatos bibliogr\'aficos contra Crossref y los registros de
los editores; asistencia con LaTeX y construcci\'on del fichero BibTeX.\\
\addlinespace
S. Hern\'andez &
Claude (Claude Code, Opus~5) &
Fase anterior: asistencia con los \emph{scripts} de auditor\'ia del corpus y
con la lectura de sus salidas. Toda magnitud atribuida a la auditor\'ia en la
Secci\'on~\ref{sec:trackanalysis} es la salida de un \emph{script} que se
ejecut\'o sobre los datos descargados, no una cifra producida por un modelo
de lenguaje.\\
\bottomrule
\end{tabularx}
\end{table}

\noindent\textbf{Verificaci\'on realizada hasta ahora.} Cada uno de los
siete trabajos citados fue recuperado y le\'ido a texto completo antes de ser
citado, y todo valor atribuido a una fuente en este documento ( tama\~nos
de conjuntos de datos, valores de m\'etricas, recuentos de estudios ) fue
transcrito de la tabla, figura o pasaje correspondiente de esa fuente. Todos
los DOI se resolvieron contra Crossref o el registro del editor, y la
bibliograf\'ia cita la versi\'on que efectivamente se ley\'o. Una entrada
merece menci\'on expl\'icita: para el marco de la fuga de datos citamos la
versi\'on de conferencia de KDD'11, que es la versi\'on obtenida y le\'ida;
existe una versi\'on de revista ampliada, que difiere en la lista de autores,
y solo sustituir\'a a la actual si se lee a trav\'es de la suscripci\'on
institucional.

\noindent\textbf{Verificaci\'on a\'un pendiente.} Una relectura manual de las
siete fuentes por parte del autor, independiente de la extracci\'on descrita
arriba, no se ha completado todav\'ia y est\'a programada antes de la
entrega. Hasta entonces, esta declaraci\'on afirma extracci\'on asistida por
m\'aquina con trazabilidad a nivel de fuente, y no confirmaci\'on humana
independiente de cada cifra. El registro de trazabilidad, afirmaci\'on,
fuente, ubicaci\'on en la fuente, se mantiene junto a esta versi\'on de
modo que cada enunciado pueda comprobarse individualmente.

\noindent\textbf{Lo que no fue generado por inteligencia artificial.} La
pregunta de investigaci\'on, las hip\'otesis de trabajo, la elecci\'on del
Track A, la decisi\'on de agrupar por episodio, la lista negra de campos
prohibidos, el dise\~no de la auditor\'ia y su ejecuci\'on, y la
interpretaci\'on de los resultados de la auditor\'ia son trabajo propio de
los autores. En este documento no se reporta ning\'un resultado experimental,
de modo que ninguna herramienta gener\'o resultado alguno.

% ===========================================================================
% CONCLUSION
% ===========================================================================
\section{Estado Actual y Pr\'oximos Pasos}
\label{sec:conclusion}

La revisi\'on establece que predecir el desenlace de una partida a partir de
un \'unico estado intermedio es un problema supervisado bien planteado en
esta familia de juegos, que las caracter\'isticas tabulares extra\'idas del
estado transportan se\~nal sustancial, y que los modelos cl\'asicos capturan
la mayor parte de ella: la mejor soluci\'on de \emph{Hearthstone} mejor\'o al
\emph{baseline} de una red de dos capas ocultas en menos de un dos por ciento
de AUC~\cite{janusz_hearthstone_2017}, y tres clasificadores convencionales
quedaron a menos de 0.016 de acierto unos de otros en
SC2EGSet~\cite{bialecki_sc2egset_2023}. Establece tambi\'en lo que la
literatura no ha resuelto. Ning\'un trabajo publicado plantea la pregunta
para el \emph{Pok\'emon TCG}. Ning\'un estudio comparable sobre estados de
juego reporta si sus probabilidades est\'an calibradas, aunque la evidencia
sistem\'atica identifica la calibraci\'on como la propiedad que falla
primero~\cite{guo_temporalshift_2021}. Y el costo de ignorar la dependencia
entre observaciones ha sido cuantificado en otros campos, un acierto de
0.91 para \emph{extra trees} que colapsa a 0.66 una vez respetada la
agrupaci\'on~\cite{karbalaie_participantaware_2026}, pero se aplica de
forma desigual dentro de los conjuntos de datos de juegos.

Este trabajo ataca esa intersecci\'on. Su objeto es el valor de un estado de
forma aislada, aprendido \emph{offline} bajo un protocolo cuyos controles de
fuga se justifican por medici\'on y no por argumento: en nuestra propia
auditor\'ia de 52\,085 episodios y 7\,038\,378 instant\'aneas, una
partici\'on aleatoria a nivel de fila situar\'ia el $99.995\%$ de los
episodios a ambos lados de la divisi\'on. La contribuci\'on a la que el
proyecto aspira no es una puntuaci\'on de titular, sino una respuesta a tres
preguntas ligadas: cu\'anto dice un estado, en qu\'e momento de la partida
empieza a decirlo, y cu\'anto sobrevive a un cambio de episodios, de mazos y
de d\'ias.

El punto de referencia declarado es la banda de $0.78$--$0.80$ de AUC del
AAIA'17 sobre \emph{Hearthstone}, adoptada como orden de magnitud y no como
banco de pruebas, con las diferencias de juego, poblaci\'on de agentes,
partici\'on, prevalencia de clases y descripci\'on del estado que impiden una
comparaci\'on estricta enunciadas en la
Secci\'on~\ref{subsec:litbaseline}. Si el modelo cl\'asico m\'as fuerte
alcanza esa banda en la ventana tard\'ia de la partida bajo nuestro protocolo
agrupado, el hallazgo sustantivo ser\'a en qu\'e punto de la partida lo hace
por primera vez.

\noindent\textbf{Pr\'oximos pasos (Etapa 2).}
\begin{itemize}
  \item Materializar la representaci\'on tabular una sola vez, en formato
        columnar, y completar sobre ella las comprobaciones consignadas como
        pendientes en la Secci\'on~\ref{sec:trackanalysis}: prevalencia de
        clases del objetivo final, ausencia de valores por columna, rangos
        imposibles, filas duplicadas y cardinalidad efectiva.
  \item Implementar los controles de fuga como pruebas ejecutables, la
        aserci\'on de lista negra y la aserci\'on de intersecci\'on vac\'ia
        entre particiones, y convertirlas en precondici\'on de cualquier
        ejecuci\'on de entrenamiento.
  \item Fijar las ventanas de turno a partir de la distribuci\'on emp\'irica
        de duraci\'on y declararlas antes de examinar la partici\'on de
        prueba.
  \item Implementar y comparar la escalera de \emph{baselines} B0 a B4 bajo
        el protocolo agrupado, reportando discriminaci\'on y calibraci\'on
        con intervalos de confianza remuestreados sobre episodios.
  \item Ejecutar los protocolos temporal y de reserva de mazos y las
        ablaciones por familia de caracter\'isticas, reportando cada uno por
        separado, y poner a prueba la hip\'otesis H3 sobre si la
        calibraci\'on se degrada antes que la discriminaci\'on.
  \item Fusionar esta revisi\'on individual con las versiones de los Ejes A y
        B producidas por los otros dos integrantes del equipo en el
        art\'iculo conjunto de la Etapa 1, resolviendo el solapamiento
        tem\'atico en lugar de concatenar las tres secciones de trabajos
        relacionados.
\end{itemize}

\bibliographystyle{IEEEtran}
\bibliography{references}

\end{document}
